% ============================================================================
% ТКАНЬ БЫТИЯ: иллюстрированная онтология Унитарного Голономного Монизма
% Популярная книга-энциклопедия: минимум математики, максимум смысла.
% Компилировать ТОЛЬКО xelatex (fontspec/polyglossia, кириллица).
% ============================================================================

\documentclass[11pt,letterpaper]{article}

% === ПАКЕТЫ (XeLaTeX) ===
\usepackage{fontspec}
\usepackage{polyglossia}
\setdefaultlanguage{russian}
\setotherlanguage{english}
\setmainfont{CMU Serif}[Extension=.otf, UprightFont=cmunrm, BoldFont=cmunbx, ItalicFont=cmunti, BoldItalicFont=cmunbi, Script=Cyrillic]
\setsansfont{CMU Sans Serif}[Extension=.otf, UprightFont=cmunss, BoldFont=cmunsx, ItalicFont=cmunsi, BoldItalicFont=cmunso, Script=Cyrillic]
\setmonofont{CMU Typewriter Text}[Extension=.otf, UprightFont=cmuntt, BoldFont=cmuntb, ItalicFont=cmunit, BoldItalicFont=cmuntx, Script=Cyrillic]
\usepackage{amsmath,amssymb}
\usepackage{tikz}
\usetikzlibrary{arrows.meta, positioning, calc, decorations.pathmorphing, shapes.geometric, backgrounds}
\usepackage{xcolor}
\usepackage{booktabs}
\usepackage{hyperref}
\usepackage{geometry}
\geometry{margin=1in}
\usepackage{enumitem}
\usepackage{caption}
\usepackage{tcolorbox}
\tcbuselibrary{skins,breakable}
\usepackage{float}
\usepackage{xspace}
\usepackage{microtype}
\usepackage{fancyhdr}
\setlength{\headheight}{14pt}
\pagestyle{fancy}
\fancyhf{}
\fancyhead[L]{\footnotesize\itshape Ткань бытия}
\fancyhead[R]{\footnotesize\thepage}
\renewcommand{\headrulewidth}{0.2pt}
\renewcommand{\footrulewidth}{0pt}

% === ЦВЕТА === (фирменная палитра: сдержанность, почти белые фоны, один акцент на голос)
\definecolor{plosblue}{HTML}{2563AE}
\definecolor{plosgray}{HTML}{566270}
\definecolor{ember}{HTML}{A8650B}      % янтарь костра
\definecolor{leaf}{HTML}{136B44}       % зелень вопросов
\definecolor{fact}{HTML}{1F4E79}       % синь точности
\definecolor{softsun}{HTML}{FCF8F2}
\definecolor{softsky}{HTML}{F4F8FB}
\definecolor{softleaf}{HTML}{F3F9F5}
\definecolor{dusk}{HTML}{5D4788}
% семь голосов (цвета измерений)
\definecolor{cA}{HTML}{C0392B} % Артикуляция
\definecolor{cS}{HTML}{CA6F1E} % Структура
\definecolor{cD}{HTML}{B7950B} % Динамика
\definecolor{cL}{HTML}{1E8449} % Логика
\definecolor{cE}{HTML}{6B4C9A} % Интериорность
\definecolor{cO}{HTML}{34495E} % Основание
\definecolor{cU}{HTML}{2E86C1} % Единство

% === ТРИ ГОЛОСА КНИГИ === (фирменный стиль: прямые углы, волосяная рамка,
%     одна весовая линейка слева, заголовок цветом рамки на светлом фоне)
\newtcolorbox{skazbox}[1][]{enhanced, colback=softsun, colframe=ember!80!black,
  coltitle=ember!70!black, colbacktitle=softsun, titlerule=0mm,
  fonttitle=\bfseries\itshape, title={#1}, boxrule=0.3pt, leftrule=2.2pt,
  arc=0pt, breakable, left=8pt, right=8pt, toptitle=3pt, bottomtitle=1pt}
\newtcolorbox{factbox}[1][]{enhanced, colback=softsky, colframe=fact,
  coltitle=fact, colbacktitle=softsky, titlerule=0mm,
  fonttitle=\bfseries, title={#1}, boxrule=0.3pt, leftrule=2.2pt,
  arc=0pt, breakable, left=8pt, right=8pt, toptitle=3pt, bottomtitle=1pt}
\newtcolorbox{askbox}[1][]{enhanced, colback=softleaf, colframe=leaf,
  coltitle=leaf!85!black, colbacktitle=softleaf, titlerule=0mm,
  fonttitle=\bfseries, title={#1}, boxrule=0.3pt, leftrule=2.2pt,
  arc=0pt, breakable, left=8pt, right=8pt, toptitle=3pt, bottomtitle=1pt}

% === КОМАНДЫ ===
\newcommand{\Gam}{\Gamma}
\newcommand{\stT}{\textbf{[Т]}\xspace}
\newcommand{\stC}{\textbf{[С]}\xspace}
\newcommand{\stP}{\textbf{[П]}\xspace}
\newcommand{\stI}{\textbf{[И]}\xspace}
\newcommand{\stH}{\textbf{[Г]}\xspace}
\newcommand{\holonsh}{\href{https://holon.sh}{holon.sh}\xspace}


% === ПОРОГ ЧАСТИ ===
\newcommand{\partline}[2]{%
  \clearpage\vspace*{4pt}\begin{center}
  {\color{fact}\rule{0.55\textwidth}{0.4pt}}\\[9pt]
  {\large\bfseries #1}\\[7pt]
  \begin{minipage}{0.82\textwidth}\centering\small\itshape\color{plosgray} #2\end{minipage}\\[9pt]
  {\color{fact}\rule{0.55\textwidth}{0.4pt}}
  \end{center}\bigskip}

% === ФИРМЕННЫЙ ЗНАК: роза семи голосов ===
% \voicerose{x}{y}{радиус}{точка}: цветная роза в (x,y)
\newcommand{\voicerose}[4]{%
  \foreach \vang/\vcol in {90/cA,141.4/cS,192.8/cD,244.3/cL,295.7/cE,347.1/cO,38.6/cU}{%
    \draw[\vcol!55, thick] (#1,#2) -- ($(#1,#2)+(\vang:#3)$);
    \fill[\vcol] ($(#1,#2)+(\vang:#3)$) circle (#4);}}
% \voiceroseneutral{x}{y}{радиус}{точка}: тихая роза для фигур населения
\newcommand{\voiceroseneutral}[4]{%
  \foreach \vang in {90,141.4,192.8,244.3,295.7,347.1,38.6}{%
    \draw[plosgray!40] (#1,#2) -- ($(#1,#2)+(\vang:#3)$);
    \fill[plosblue!65] ($(#1,#2)+(\vang:#3)$) circle (#4);}}

\hypersetup{colorlinks=true, linkcolor=fact, urlcolor=plosblue, citecolor=plosgray,
  pdftitle={Ткань бытия: иллюстрированная онтология Унитарного Голономного Монизма},
  pdfauthor={Корпус УГМ (holon.sh)},
  pdfsubject={Популярная онтология УГМ: математическая оптика и когерентная кибернетика},
  pdfkeywords={УГМ, онтология, сознание, квалиа, когерентность, плоскость Фано, октонионы}}
\captionsetup{font=small, labelfont={bf,color=fact}}

\begin{document}

% ============================= ТИТУЛ =============================
\begin{titlepage}
\centering
\vspace*{28pt}
{\Huge\bfseries Ткань бытия}\\[10pt]
{\LARGE Иллюстрированная онтология\\[3pt] Унитарного Голономного Монизма}\\[26pt]
% --- титульная мандала семи голосов ---
\begin{tikzpicture}[scale=1.25]
  \foreach \ang/\col/\lab in {90/cA/A, 141.4/cS/S, 192.8/cD/D, 244.3/cL/L, 295.7/cE/E, 347.1/cO/O, 38.6/cU/U}{
    \draw[\col!45, line width=1pt] (0,0) -- (\ang:2.1);
    \fill[\col] (\ang:2.1) circle (0.19);
    \node[font=\small\bfseries, text=white] at (\ang:2.1) {\lab};
  }
  \foreach \a in {90, 141.4, 192.8, 244.3, 295.7, 347.1, 38.6}{
    \foreach \b in {90, 141.4, 192.8, 244.3, 295.7, 347.1, 38.6}{
      \draw[plosgray!25, line width=0.3pt] (\a:2.1) -- (\b:2.1);
    }}
  \fill[white] (0,0) circle (0.55);
  \draw[plosgray!60] (0,0) circle (0.55);
  \node[font=\footnotesize\itshape, text=plosgray] at (0,0.12) {один};
  \node[font=\footnotesize\itshape, text=plosgray] at (0,-0.16) {аккорд};
\end{tikzpicture}\\[26pt]
{\large Аксиоматическая оптика для исследований и практики ---\\
для всякой области, где что-то живёт, мыслит или строится}\\[18pt]
{\small\color{plosgray} По теории УГМ (\holonsh) --- с честными пометками,\\
что доказано, что предположено и что пока лишь красиво}\\[10pt]
{\small\color{plosgray} Издание первое --- август 2026}
\vfill
\begin{skazbox}[Как читать эту книгу]
В книге три голоса. \textcolor{ember}{\textbf{Янтарные врезки}} --- сказка у
костра: образы, которые можно читать ребёнку. \textcolor{fact}{\textbf{Синие}}
--- «если точно»: что теория утверждает на самом деле, без украшений.
\textcolor{leaf}{\textbf{Зелёные}} --- вопросы, о которых стоит поспорить за
ужином. Основной текст --- для всех остальных: инженеров, врачей, философов,
психологов --- взрослых людей, которым нужна суть без формул. Формул здесь
почти нет; там, где число важно, оно появляется как герой рассказа, а не как
уравнение.
\end{skazbox}
\end{titlepage}

\tableofcontents
\newpage

% ============================= ЛЕГЕНДА ЧЕСТНОСТИ =============================
\section*{Пять пометок честности}
\addcontentsline{toc}{section}{Пять пометок честности}

Наука отличается от мифа не отсутствием красоты, а привычкой помечать,
что именно она знает. Всюду в книге встретятся пять пометок:

\begin{center}
\begin{tabular}{@{}cl@{}}
\toprule
\stT & \textbf{доказано}: теорема внутри теории; спорить можно только с аксиомами \\
\stC & \textbf{доказано при условии}: верно, если верна названная посылка \\
\stH & \textbf{гипотеза}: правдоподобно, ищется доказательство \\
\stP & \textbf{постулат}: принято как исходная точка, не выведено \\
\stI & \textbf{интерпретация}: прочтение, а не факт; красиво, но не принуждает \\
\bottomrule
\end{tabular}
\end{center}

Когда в книге сказано «мир устроен так», рядом всегда стоит пометка.
Это и есть главное отличие этой энциклопедии от священных книг: она
говорит, \emph{как проверить, что она ошибается} (глава~\ref{sec:check}).

За пятью пометками стоит один закон. Утверждение \textbf{зрело},
когда держит три обязательства разом: оно \emph{показывает, как
построено} --- конкретная конструкция, а не обещание существования;
оно \emph{отдаётся механической проверке} --- его может проверить
машина, без апелляций к авторитету и интуиции; и оно
\emph{работает} --- есть действующее представление: программа,
процедура, прибор. Пометки --- сжатые отпечатки этих трёх: \stT{}
--- все три разом; \stC{} --- все три при названных посылках;
\stH{} --- формулировка с планом, но пока без доказательства;
\stP{} --- несущее допущение, объявленное как допущение; \stI{} ---
прочтение, которое книга держит наперечёт и на виду.

Зачем такая строгость? Потому что знание, рассчитанное пережить
своих создателей, встречает три эрозии: дрейфуют слова, сменяются
инструменты, смешиваются регистры --- утверждения, надежды и
риторика становятся неотличимы. Три обязательства --- защита от
всех трёх сразу: конструкции переживают перевод, механические
проверки переживают смену инструментов, помеченные регистры
переживают дрейф стиля. Отсюда одно домашнее правило, ощутимое на
каждой странице: \textbf{в этой книге нет непомеченных
утверждений}. Что не смогло понести свою пометку --- было
достроено, пока не смогло, честно понижено --- или удалено. Вы
читаете то, что выжило.

% ============================= ЧТО ПЕРЕД ВАМИ =============================
\section*{Что перед вами: оптика, а не доктрина}
\addcontentsline{toc}{section}{Что перед вами: оптика, а не доктрина}

УГМ полезнее всего держать в руках не как «ещё одну картину мира», а
как \textbf{математическую оптику}: один жёстко выведенный объект ---
семь голосов, их дружбы, единый закон изменения, --- сквозь который
по очереди наводится на резкость всё остальное: частицы и
пространство, жизнь и смерть, чувство и слово, семья и цивилизация.
От доктрины оптику отличают два свойства. Она \emph{выводится}, а не
проповедуется: почему у линзы ровно семь элементов --- теорема, а не
вкус (глава~\ref{sec:seven}). И она \emph{наводится}: каждое понятие
этой книги --- собранность, запас, сцепленность, краска, стресс
голоса --- не метафора, а вычислимая величина с порогом; у мира,
прочитанного этой оптикой, появляются приборные панели
(глава~\ref{sec:health}) и паспорта (глава~\ref{sec:qualia}).

Из оптики вырастает переосмысление жизни, которое корпус называет
\textbf{когерентной кибернетикой}: жить --- значит удерживать
согласие семи голосов на подпитке (глава~\ref{sec:life});
заботиться --- значит восстанавливать поток, не гася чужих этажей
(глава~\ref{sec:health}); этика перестаёт быть списком запретов и
становится инженерией согласий --- с теоремами там, где прежде
стояли благие пожелания: связь прибавляет бытия
(глава~\ref{sec:others}), сверхразум растёт лесом, а не башней
(глава~\ref{sec:society}).

И последнее, что отличает рабочую оптику от красивой гипотезы: ею
уже работают. Вердикты «машинно проверено» этой книги --- испытания
одной и той же живой машинерии: на этой самой семёрке выращена
действующая когнитивная архитектура, и главы о золотых путях, языке
и красках пересказывают её измерения, а не мысленные опыты. Её
устройство --- не предмет этой книги; предмет --- сам факт: оптика
достаточно строга, чтобы нести вес машин. Это не проверка
предсказаний о природе --- те собраны в реестре
(глава~\ref{sec:check}); это проверка пригодности словаря. Для
учёного и инженера в этом и главный подарок теории: язык, на
котором о живом можно не только говорить --- на нём можно строить.

Взятое вместе, это делает УГМ не мировоззрением, которое принимают,
а \textbf{аксиоматической исследовательской программой}, в которой
работают: жёсткое ядро --- одна аксиома связи, семь вынужденных
измерений, теоремы, проверенные машиной, константы, пересчитываемые
на ноутбуке, --- и растущий вокруг него пояс приложений, каждое со
своей пометкой и своим способом провалиться. Книга построена как
одна дорога через эту программу, в четыре движения. Сначала
\emph{мир}: что такое ткань и почему в её алфавите ровно семь букв.
Затем \emph{внутри}: как на этой ткани загораются жизнь, окно
сознания, этажи рефлексии и краски переживания. Затем \emph{узор в
действии}: свобода, смелые пути, язык, память --- механика
действующей жизни. И наконец \emph{вместе --- и что теперь}:
другие, общества, забота, смерть, наблюдатель --- и в конце не
итог, а рабочая программа: что эта оптика кладёт на стол физику,
врачу, инженеру разума, учителю --- и где она сейчас уязвима. Какова
бы ни была ваша область, задуманный опыт этой книги один: где-то в
этих главах ваш собственный самый трудный вопрос должен вдруг
появиться --- пересказанным на языке, где его можно вычислить,
проверить или честно пометить открытым.

\partline{ЧАСТЬ ПЕРВАЯ \;\textperiodcentered\; МИР}{Из чего соткано всё: аккорд семи голосов,
вынужденность семёрки, рождение времени, начала, пространства и вещей. Здесь строится сцена ---
и на ней пока никого нет.}
\addcontentsline{toc}{section}{ЧАСТЬ ПЕРВАЯ · МИР}

% ============================= 1. ИЗ ЧЕГО ВСЁ =============================
\section{Из чего сделано всё}
\label{sec:fabric}

Спросите физика прошлого века, из чего сделан мир, --- он ответит:
из маленьких кирпичиков, частиц. Спросите УГМ --- ответ другой и
неожиданно музыкальный: \emph{мир сделан из согласия}. Не из вещей,
а из способа, которым вещи держатся друг за друга.

Представьте семь голосов, поющих один аккорд. Уберите певцов, оставьте
\emph{сам аккорд} --- соотношение высот, взаимную настройку, дрожание
согласных звуков. В УГМ первично именно это: не «что звучит», а «как
согласовано». Математики называют такой объект матрицей плотности и пишут
$\Gam$: семь её диагональных ячеек --- силы семи голосов, а двадцать
одна ячейка над диагональю --- их попарные дружбы (когерентности).
Мы будем называть его \textbf{аккордом мира}.

У аккорда семь голосов --- семь измерений всякого существования:

\begin{center}
\begin{tikzpicture}[font=\small]
  \foreach \x/\col/\ltr/\name/\verb in {0/cA/A/Артикуляция/различать, 2.3/cS/S/Структура/{держать форму},
    4.6/cD/D/Динамика/двигаться, 6.9/cL/L/Логика/согласовывать, 9.2/cE/E/Интериорность/переживать,
    11.5/cO/O/Основание/укоренять, 13.8/cU/U/Единство/{собирать в одно}}{
    \fill[\col!85] (\x,0) circle (0.42);
    \node[font=\footnotesize\bfseries, text=white] at (\x,0) {\ltr};
    \node[font=\footnotesize\bfseries] at (\x,-0.85) {\name};
    \node[font=\scriptsize\itshape, text=plosgray] at (\x,-1.25) {\verb};
  }
\end{tikzpicture}
\end{center}

Различать, держать форму, двигаться, согласовывать, переживать,
укоренять, собирать в одно. За каждым словом --- строгая величина
корпуса теории: Артикуляция (A) --- активность различения; Структура
(S) --- устойчивость формы; Динамика (D) --- активность процесса;
Логика (L) --- внутренняя согласованность; Интериорность (E) ---
насыщенность внутренней стороны; Основание (O) --- связь с истоком;
Единство (U) --- интеграция в целое. Это не список органов и не
поэтическая вольность:
дальше мы увидим, что семь --- не выбор автора теории, а вынужденный
ответ математики (глава~\ref{sec:seven}).

\begin{skazbox}[Сказка у костра --- про оркестр без музыкантов]
Жил-был оркестр, в котором не было музыкантов. Были только
\emph{дружбы между нотами}: эта нота тянется к той, та спорит с
третьей, а все вместе они держат один долгий аккорд. Когда дружбы
крепкие --- аккорд звонкий, и в мире появляется вещь: камень, кошка,
ты. Когда дружбы рвутся --- аккорд рассыпается в шорох, и вещь
растворяется. Мир --- это не склад кирпичей. Мир --- это оркестр
дружб.
\end{skazbox}

\begin{factbox}[Если точно]
Первичный объект УГМ --- не частица и не поле, а состояние
семи взаимосвязанных измерений (матрица когерентности $\Gam$ размера
$7{\times}7$) и единый закон его изменения. Всё остальное ---
частицы, пространство, время, жизнь, сознание --- \emph{выводится}
как узоры этого объекта. Статус: сама семёрка и закон --- теоремы
внутри аксиоматики \stT; выбор аксиом --- постулат \stP.
\end{factbox}

\begin{askbox}[Подумай сам]
Что «более настоящее»: скрипка --- или строй, в котором она звучит с
другими? Если бы ты был не телом, а согласием своих частей, --- что
изменилось бы в твоём отношении к себе?
\end{askbox}

% ============================= 2. ПОЧЕМУ СЕМЬ =============================
\section{Почему именно семь}
\label{sec:seven}

Число семь в этой книге --- не магия и не вкус автора. К нему ведут
четыре независимые дороги, и все четыре упираются в одно и то же.

\begin{figure}[H]
\centering
\begin{tikzpicture}[font=\small, >=Stealth]
  \node[draw=cL!70, fill=cL!8, rounded corners, text width=44mm, align=center, inner sep=5pt] (r1) at (-4.1,2.2)
    {\textbf{Дорога чисел}\\ \footnotesize системы, где можно делить, кончаются на октонионах; их «рабочая часть» семимерна};
  \node[draw=cE!70, fill=cE!8, rounded corners, text width=44mm, align=center, inner sep=5pt] (r2) at (4.1,2.2)
    {\textbf{Дорога симметрии}\\ \footnotesize у октонионов есть собственная группа поворотов ($G_2$); её сцена --- семь «мнимых» направлений октонионов};
  \node[draw=cO!70, fill=cO!8, rounded corners, text width=44mm, align=center, inner sep=5pt] (r4) at (-4.1,-2.2)
    {\textbf{Дорога зеркала}\\ \footnotesize наименьший алфавит, в котором любая одиночная поломка сама себя называет, --- семибуквенный};
  \node[draw=cS!70, fill=cS!10, rounded corners, text width=44mm, align=center, inner sep=5pt] (r3) at (4.1,-2.2)
    {\textbf{Дорога узора}\\ \footnotesize самая маленькая «плоскость из точек и прямых» (плоскость Фано) --- ровно 7 точек и 7 прямых};
  \node[circle, draw=dusk, fill=dusk!12, very thick, minimum size=17mm, font=\Large\bfseries] (seven) at (0,0) {7};
  \draw[->, thick, plosgray] (r1)--(seven);
  \draw[->, thick, plosgray] (r2)--(seven);
  \draw[->, thick, plosgray] (r3)--(seven);
  \draw[->, thick, plosgray] (r4)--(seven);
  \draw[<->, densely dashed, leaf, thick] (r1.west) to[bend right=40]
    node[pos=0.5, left=3pt, font=\scriptsize, text=leaf, align=center]{одна\\ и та же\\ теорема\\ \stT} (r4.west);
\end{tikzpicture}
\caption{Четыре дороги к семёрке. Особенно красиво, что дорога чисел и
дорога зеркала --- не совпадение, а доказанное тождество \stT: «последняя
система, где делится» и «последний алфавит, который сам себя проверяет» ---
это \emph{одно и то же} свойство, увиденное с двух сторон.}
\label{fig:seven}
\end{figure}

Семь голосов соединены не как попало: их дружбы разложены по семи
«прямым» плоскости Фано --- по три голоса на прямой, каждая пара
голосов ровно на одной прямой. Это и есть \textbf{алфавит связи} мира:
кто с кем образует троицу. Та же схема, независимо переоткрытая людьми
в теории кодов, умеет находить любую одиночную ошибку --- поэтому
мир, сотканный по Фано, \emph{умеет замечать свои поломки}. Запомните
это: в главе о жизни оно вернётся.

\begin{skazbox}[Сказка у костра --- про семь гномов и карту дружб]
У семи гномов была карта: семь домиков и семь тропинок, на каждой
тропинке --- по три домика. Гномы заметили странное: если один из них
заболевал, тропинки \emph{сами} показывали, кто именно, --- по тому,
какие из них начинали скрипеть. «Наша карта умеет жаловаться», ---
сказали гномы. А мудрый гном добавил: «Потому нас и семь. Будь нас
шесть или восемь --- карта бы молчала или путалась».
\end{skazbox}

\begin{askbox}[Подумай сам]
Почему «самопроверяющийся алфавит» --- полезное свойство для всего
живого? Придумай, что было бы с существом, чьи поломки никак себя не
называют.
\end{askbox}

% ============================= 2b. ЧИСЛА =============================
\section{Почему числа именно такие: лестница из четырёх миров}
\label{sec:numbers}

Прежде чем идти дальше, стоит спросить о самом инструменте. Мир в этой
книге описывается числами --- обычными, комплексными, и той самой
семёркой мнимых направлений. Но откуда взялись \emph{сами числа}?
Почему у мироздания именно такой счётный инструмент, а не другой?
УГМ отвечает: инструмент не выбран --- он \emph{вырос}, и лестница
его роста коротка и обрывается сама.

Есть простое правило --- назовём его правилом удвоения: если у тебя
есть мир чисел, где можно складывать, умножать и делить, из него можно
построить мир вдвое больше, склеив две копии. Так из обычной прямой
$\mathbb{R}$ (одно измерение) вырастает плоскость $\mathbb{C}$ (два),
из неё --- четырёхмерный мир кватернионов $\mathbb{H}$, из него ---
восьмимерный мир октонионов $\mathbb{O}$, чья мнимая часть и есть наша
семёрка. Лестница $1 \to 2 \to 4 \to 8$ --- и на каждой ступени
приходится платить: сначала теряется порядок (нельзя сказать, какое
комплексное число «больше»), потом перестановочность умножения, потом
его группировка. А на следующем шаге --- шестнадцатимерные седенионы
--- происходит непоправимое: появляются \emph{делители нуля}, пары
ненулевых чисел, дающие в произведении ноль. Число, помноженное на
число, --- и вдруг ничто. Мир с такой трещиной не может нести
различий: в нём «что-то» умеет бесследно исчезать. Лестница
обрывается на восьми --- это теорема \stT, доказанная ещё в XIX веке
и перепроверенная в корпусе теории машинным перебором.

Вот почему семёрка глав~\ref{sec:seven} --- не одна из возможных, а
\emph{последняя возможная}: мнимая часть последнего живого мира чисел.

\begin{figure}[H]
\centering
\begin{tikzpicture}[font=\small, >=Stealth]
  % лестница
  \foreach \i/\x/\lab/\dim in {0/-5.2/$\mathbb{R}$/1, 1/-3.2/$\mathbb{C}$/2, 2/-1.2/$\mathbb{H}$/4, 3/0.8/$\mathbb{O}$/8}{
    \draw[fact!75, fill=softsky] (\x,\i*0.55) rectangle (\x+1.5,\i*0.55+0.5);
    \node at (\x+0.75,\i*0.55+0.25) {\lab\ \scriptsize(\dim)};
  }
  \foreach \i/\x in {0/-3.62, 1/-1.62, 2/0.38}{
    \draw[->, thick, plosgray] (\x,\i*0.55+0.28) -- (\x+0.38,\i*0.55+0.72);
  }
  % обрыв: седенионы
  \draw[->, thick, densely dashed, plosgray] (2.35,2.0) -- (2.95,2.32);
  \draw[plosgray!60, fill=plosgray!12, dashed] (3.0,2.2) rectangle (4.7,2.7);
  \node[text=plosgray] at (3.85,2.45) {16?};
  \draw[ember, very thick] (3.25,2.15) -- (4.45,2.75);
  \draw[ember, very thick] (3.25,2.75) -- (4.45,2.15);
  \node[font=\scriptsize, text=ember, align=center] at (3.85,3.25)
    {делители нуля:\\ число $\times$ число $=$ ничто \stT};
  \node[font=\scriptsize, text=plosgray, align=center] at (-2.6,3.25)
    {правило удвоения: $1 \to 2 \to 4 \to 8$ --- и обрыв};
  % уроборос: тело с разрывом-дырой, голова тянется к хвосту
  \draw[dusk, very thick, -{Stealth[length=7pt]}] ($(7.0,1.2)+(150:1.05)$) arc (150:425:1.05);
  \fill[dusk] ($(7.0,1.2)+(150:1.05)$) circle (0.06);
  \draw[ember, very thick, densely dotted] ($(7.0,1.2)+(150:1.05)$) arc (150:65:1.05);
  \node[font=\scriptsize, text=ember] at ($(7.0,1.2)+(107:1.5)$) {дыра};
  \node[font=\scriptsize, text=dusk, align=center] at (7.0,-0.45)
    {змея, кусающая хвост:\\ если в прямой есть дыра ---\\ хвост не поймать};
\end{tikzpicture}
\caption{Слева: лестница миров чисел. Каждая ступень --- удвоение с
платой; на шестнадцати появляются делители нуля --- трещина, сквозь
которую исчезают различия, --- и лестница обрывается \stT. Справа:
почему прямая чисел \emph{сплошная} --- замыкание самомодели (змея
обязана дотянуться до собственного хвоста) невозможно над дырявой
прямой; непрерывность --- не допущение, а следствие уробороса \stT.}
\label{fig:numbers}
\end{figure}

И второе, тихое чудо этой главы. Почему прямая чисел \emph{сплошная},
без дыр? В корпусе доказано \stT{}: это следует из самого древнего
образа теории --- змеи, кусающей свой хвост. Мир, который моделирует
сам себя (глава~\ref{sec:observer}), обязан уметь \emph{дотянуться до
себя}: у непрерывного самоотображения должна существовать неподвижная
точка. Над дырявой прямой --- как над одними дробями --- змея
буквально перепрыгивает собственный хвост сквозь дыру, и замыкание
срывается. Требование «змея всегда ловит хвост» \emph{вынуждает}
сплошную прямую --- а вместе с ней и непрерывное время: поток закона
корректен, потому что числа полны. Время течёт гладко не по
привычке мира, а потому, что мир себя держит.

\begin{skazbox}[Сказка у костра --- про кузнеца и четыре слитка]
Кузнец ковал миры для счёта. Первый слиток --- прут: прямой, честный,
всё на нём по порядку. Второй --- лист: на нём уже можно кружить.
Третий --- брусок: в нём можно вертеть. Четвёртый --- восьмигранный
колокол: он один умел звенеть семью голосами сразу. Замахнулся кузнец
на пятый --- а слиток треснул: ударишь по нему, а звона нет ---
глотает сам себя. «Дальше не куётся», --- сказал кузнец. И оставил
колокол. С тех пор всё сущее звенит в семь голосов.
\end{skazbox}

\begin{askbox}[Подумай сам]
Правило удвоения даёт $1, 2, 4, 8$ --- и смерть на $16$. Найди в
своей жизни лестницу, которая растёт удвоением и где-то обрывается
сама: команда, которая ещё дружит; список дел, который ещё держится в
голове. Что играет роль «делителя нуля» --- признака, что следующая
ступень уже не живая?
\end{askbox}

% ============================= 3. ВРЕМЯ =============================
\section{Что такое время}
\label{sec:time}

В обыденной картине время --- река: она течёт сама по себе, а мы
плывём. УГМ говорит строже и интереснее: \emph{реки нет}. Полное
состояние мира вообще не меняется --- меняются лишь его части
\emph{друг относительно друга}. Время --- это не сцена, а танец.

Представьте двух танцоров без всякой музыки. Ни один не движется «во
времени» --- но каждый движется \emph{в такт другому}. Если спросить
первого, «который час», он ответит, взглянув на позу второго. В УГМ
роль такого партнёра играет один из семи голосов --- Основание (O):
его размеренное «дыхание» служит часами для всех остальных. Часы не
снаружи мира; часы --- один из его голосов.

\begin{figure}[H]
\centering
\begin{tikzpicture}[font=\small, >=Stealth]
  \draw[plosgray!50, thick] (0,0) circle (1.5);
  \foreach \a in {0,30,...,330}{\draw[plosgray!50] (\a:1.4)--(\a:1.5);}
  \fill[cO] (90:1.05) circle (0.09);
  \node[font=\scriptsize, text=cO] at (90:1.9) {голос O --- «маятник»};
  \draw[->, thick, cO] (0,0)--(90:1.0);
  \draw[->, thick, cE!80] (0,0)--(35:0.8);
  \node[font=\scriptsize, text=cE, anchor=west] at (6:1.62) {остальные голоса};
  \node[font=\scriptsize, text=plosgray, align=center] at (0,-2.15)
    {никакого «внешнего» циферблата:\\ стрелки читают время друг по другу};
\end{tikzpicture}
\caption{Время по УГМ: не река, в которой плывёт мир, а согласованность
частей мира между собой. Вопрос «что было, когда времени ещё не было»
из осмысленного превращается в грамматическую ошибку --- как «что
севернее Северного полюса».}
\label{fig:time}
\end{figure}

\begin{factbox}[Если точно]
Полное состояние вселенной стационарно; внутреннее время возникает
относительно «часового» подсистемного голоса O (механизм
Пейджа--Вуттерса) \stT. Следствие, важное для следующей главы: пока
голос O ничем не выделен среди семи, слова «раньше» и «позже» не
определены вовсе.
\end{factbox}

% ============================= 4. ИСТОЧНИК =============================
\section{Начало, у которого не было «до»}
\label{sec:source}

Теперь мы готовы к главному вопросу всякой онтологии: \emph{как всё
началось?} И к честному ответу УГМ, у которого три слоя.

\textbf{Слой первый --- что было «в начале» \stP.} Теория постулирует
исходное состояние, которое называет \textbf{Источником}: все семь
голосов слиты в одну совершенную ноту, спетую с равной силой. Это
самое согласованное состояние, какое вообще возможно, --- и при этом
самое \emph{неразличённое}: в нём ничего ни от чего не отличается.
Ни один голос не выделен --- а значит (глава~\ref{sec:time}), нет и
часов, нет «раньше», нет «до». Источник не «был давным-давно» ---
он вне времени, как обложка вне страниц книги.

Красивая строгость: сама форма Источника --- не произвол. Доказано
\stT, что существует ровно \emph{одно} состояние, в котором ни один
голос не выделен, и ровно \emph{одно} максимально согласованное, ---
и это одно и то же состояние. Источник единствен дважды.

\textbf{Слой второй --- почему всё началось \stT.} Вот теорема,
достойная того, чтобы рассказывать её детям: \emph{совершенство
неустойчиво}. Источник --- не точка покоя. Закон мира, приложенный к
идеально слитному аккорду, немедленно даёт ненулевой сдвиг: аккорд
начинает расслаиваться. Более того, в законе спрятана петля усиления:
чуть только голос Интериорности (E) и голос Основания (O)
отличились от прочих --- их отличие начинает питать само себя. Симметрия ломается
не случайно и не по чьей-то воле, а \emph{неизбежно}, как карандаш,
поставленный на остриё, неизбежно падает --- с той разницей, что здесь
доказано и то, что «остриё» вообще не может быть точкой равновесия.

\textbf{Слой третий --- куда всё пошло \stT.} Машинная проверка
закона (его буквально интегрируют на компьютере) показывает: из
Источника мир скатывается не куда попало, а в \emph{живое} состояние
--- в то самое окно жизнеспособности, о котором вся глава~\ref{sec:life}.
Расслоение рождает различия, различия рождают часы, часы рождают
историю. По оценкам теории, первые мгновения этой истории заняли
порядка семи планковских времён --- около $4\cdot10^{-43}$ секунды \stC.

\begin{figure}[H]
\centering
\begin{tikzpicture}[font=\small, >=Stealth]
  % Источник
  \fill[dusk!25] (-5.2,0) circle (0.85);
  \draw[dusk, thick] (-5.2,0) circle (0.85);
  \node[font=\scriptsize\bfseries, text=dusk] at (-5.2,1.25) {Источник};
  \node[font=\scriptsize, text=dusk!80!black, align=center] at (-5.2,-1.3)
    {одна нота,\\ вне времени};
  % стрелка нестабильности
  \draw[->, very thick, ember] (-4.2,0) -- (-2.7,0)
    node[midway, above, font=\scriptsize, text=ember, align=center]{неустойчив\\ \stT};
  % расслоение
  \voicerose{-1.2}{0}{1.15}{0.1}
  \node[font=\scriptsize\bfseries] at (-1.2,1.6) {семь голосов};
  \node[font=\scriptsize, text=plosgray, align=center] at (-1.2,-1.62) {различие --- часы ---\\ история};
  % стрелка в окно
  \draw[->, very thick, leaf] (0.3,0) -- (1.8,0)
    node[midway, above, font=\scriptsize, text=leaf]{само};
  % окно жизни
  \draw[leaf!70, very thick] (2.6,-1.0) rectangle (5.6,1.0);
  \fill[leaf!10] (2.6,-1.0) rectangle (5.6,1.0);
  \node[font=\scriptsize\bfseries, text=leaf!50!black, align=center] at (4.1,0.32)
    {окно\\ жизнеспособности};
  \node[font=\scriptsize, text=plosgray] at (4.1,-0.55) {сюда --- в главу~\ref{sec:life}};
\end{tikzpicture}
\caption{Космогенезис по УГМ: Источник (постулат \stP) неустойчив
(теорема \stT), расслаивается в семь голосов с часами и историей и
\emph{сам собой} скатывается в окно жизни (машинно проверено).
Вопрос «что было до?» растворён: «до» --- слово, которое рождается
на втором шаге этой картинки.}
\label{fig:source}
\end{figure}

\begin{skazbox}[Сказка у костра --- про ноту, которой стало тесно]
Сначала была одна нота. Такая ровная, что в ней не было ни верха, ни
низа, ни вчера, ни завтра. Но у совершенной ноты есть тайна: ей не на
что опереться. Она дрогнула --- и из дрожи родились семь голосов. Один
стал считать удары --- и появилось «раньше» и «потом». Другой стал
слушать себя --- и появилось «внутри». Нота не сломалась. Нота
\emph{распустилась}, как бутон, который слишком совершенен, чтобы не
цвести.
\end{skazbox}

\begin{factbox}[Честный остаток]
Что здесь доказано, а что нет. Неустойчивость Источника, петля
усиления, скат в живое окно --- теоремы и машинные проверки \stT.
Сам Источник --- постулат \stP{}: «почему исходное состояние именно
таково» --- открытый вопрос. Ответ «ничто неустойчиво, потому что
несамосогласованно» --- философская позиция \stI, а не теорема, и
книга обязана это сказать вслух.
\end{factbox}

\begin{askbox}[Подумай сам]
«Совершенство неустойчиво» --- где ты встречал это в жизни? Пустой
лист, идеальная тишина, «идеальные» отношения без единого спора...
Что общего у этих неустойчивостей с Источником?
\end{askbox}

% ============================= 4b. СЦЕНА =============================
\section{Сцена: почему пространство трёхмерно}
\label{sec:space}

Мы привыкли думать: вещи находятся \emph{в} пространстве, как мебель в
комнате. УГМ переворачивает картину: \textbf{пространство находится в
вещах} --- это не коробка, в которую положили мир, а способ, которым
часть голосов держит друг друга.

Вспомните расслоение Источника (глава~\ref{sec:source}). Когда голос
Основания (O) стал часами, от хора всех поворотов мира остаётся
подхор, \emph{щадящий часы}, --- и у его наименьшей сцены ровно
\textbf{три} независимых направления: это теорема \stT, чистая
математика семёрки. Вот и вся тайна длины, ширины и высоты. Шесть
нечасовых голосов при этом делятся на две тройки: наружную, чьи
направления развёрнуты в привычную даль, и внутреннюю, чьи
направления свёрнуты в микроскопические кольца и ощущаются не как
«места», а как внутренние свойства частиц. Какие именно голоса
смотрят наружу --- теория указывает на Артикуляцию, Структуру и
Динамику, а внутрь --- на Логику, Интериорность и Единство, --- это
отождествление при одной физической гипотезе \stC; само число осей
от неё не зависит.

Больше того: у сцены есть и стрела. Одно временно́е направление и три
пространственных --- знаменитая подпись $(1{,}3)$ нашего мира ---
выводится, а не постулируется: время наследует знак от часового
голоса O, пространство --- от наружной тройки \stT{} (последний шаг ---
при одном техническом условии физической честности \stC). Даже то,
что у мира есть «вперёд» и «вширь», --- следствие всё той же семёрки.

\begin{figure}[H]
\centering
\begin{tikzpicture}[font=\small, >=Stealth]
  % центральный часовой узел
  \fill[cO] (0,0) circle (0.3);
  \node[font=\scriptsize\bfseries, text=cO] at (0,-0.72) {O --- часы};
  % наружная тройка = оси сцены
  \draw[->, very thick, cA] (0,0) -- (2.6,0.6) node[right, font=\scriptsize, text=cA]{ось 1};
  \draw[->, very thick, cS] (0,0) -- (2.2,-1.3) node[right, font=\scriptsize, text=cS]{ось 2};
  \draw[->, very thick, cD] (0,0) -- (0.4,2.0) node[above, font=\scriptsize, text=cD]{ось 3};
  \node[font=\scriptsize, text=plosgray, align=center] at (3.6,1.75)
    {наружная тройка: три оси \stT\\ какие голоса --- \stC};
  % внутренняя тройка
  \draw[->, densely dashed, cL!70] (0,0) -- (-2.3,1.0);
  \draw[->, densely dashed, cE!70] (0,0) -- (-2.5,-0.2);
  \draw[->, densely dashed, cU!70] (0,0) -- (-1.8,-1.5);
  \node[font=\scriptsize, text=plosgray, align=center] at (-2.5,-1.95)
    {внутренняя тройка:\\ изнанка сцены --- не место, а нутро};
\end{tikzpicture}
\caption{Откуда сцена: часовой голос O задаёт «когда»; повороты,
щадящие часы, дают ровно три независимых направления --- оси «где»
\stT. Наружная тройка голосов развёрнута в даль, внутренняя свёрнута
внутрь вещей (какие голоса куда --- \stC). Пространство --- не вместилище
вещей, а рисунок их согласий; далёкая звезда не «лежит в другом месте
коробки» --- она далеко \emph{по карте согласий}.}
\label{fig:stage}
\end{figure}

\begin{skazbox}[Сказка у костра --- про дом из хороводов]
Дети спросили строителя: «где ты берёшь комнаты?» Строитель ответил:
«я не беру комнаты. Я ставлю три хоровода: один кружит вправо-влево,
другой --- вперёд-назад, третий --- вверх-вниз. Там, где три хоровода
держатся за руки, само собой получается место». Так и мир: сначала
хороводы, потом комнаты. А не наоборот.
\end{skazbox}

\begin{askbox}[Подумай сам]
Попробуй объяснить, где находится «далеко», не показывая пальцем:
только через то, \emph{насколько слабо} вы с этим местом связаны.
Получилось? Ты только что пересказал определение расстояния по УГМ.
\end{askbox}

% ============================= 5. ВЕЩИ И ЧАСТИЦЫ =============================
\section{Что такое вещи (и есть ли самая маленькая)}
\label{sec:things}

Если мир --- аккорд, то что такое электрон? Ответ УГМ: \textbf{узор}.
Частица --- не крупинка вещества, а устойчивый завиток в ткани
согласий, как узел на верёвке или воронка в реке. Воронку нельзя
вынуть из реки и положить в карман --- но она совершенно реальна:
у неё есть место, форма, энергия, и она может столкнуться с другой
воронкой.

Отсюда честный ответ на вопрос «какая частица самая маленькая»:

\begin{itemize}[leftmargin=1.6em]
  \item \textbf{Самой маленькой крупинки нет вовсе} --- на дне не
        песчинки, а сам аккорд: семь голосов и их связи. Спрашивать
        «из каких частиц сделан аккорд» --- как спрашивать, из каких
        волн сделана вода.
  \item \textbf{Самый лёгкий из массивных узоров} --- лёгчайшее нейтрино:
        теория предсказывает ему массу не больше $0{,}004$~эВ ---
        примерно в сто миллионов раз легче электрона \stC. Это почти
        неощутимая рябь на ткани мира.
  \item \textbf{Самый неожиданный предсказанный узор} --- «O-квант»:
        тяжёлое возбуждение голоса Основания. Он невидим для света,
        почти ни с чем не взаимодействует и \emph{не может распасться}:
        особая чётность голоса O запрещает ему исчезать \stT. Теория
        предлагает его на роль тёмной материи --- невидимого клея
        галактик \stC.
\end{itemize}

\begin{figure}[H]
\centering
\begin{tikzpicture}[font=\small, >=Stealth]
  % тихий омут под вихрем
  \fill[plosblue!10] (0,0) circle (0.70);
  % дальние струи: ровная рябь
  \foreach \s in {1,-1}{
    \draw[plosblue!45, thick, decorate, decoration={snake, amplitude=1.1pt, segment length=15pt}]
      (-5.6,\s*0.85) -- (5.6,\s*0.85);
    % ближние струи огибают омут
    \draw[plosblue!45, thick]
      (-5.6,\s*0.35) .. controls (-2.6,\s*0.35) and (-1.6,\s*0.42) .. (-0.85,\s*0.66)
      .. controls (-0.1,\s*0.9) and (0.35,\s*0.86) .. (0.95,\s*0.62)
      .. controls (1.7,\s*0.38) and (2.6,\s*0.35) .. (5.6,\s*0.35);
  }
  % осевая струя втягивается слева и рождается справа --- по ходу вращения
  \draw[plosblue!45, thick] (-5.6,0) .. controls (-2.6,0) and (-1.5,-0.02) .. (-0.78,-0.18);
  \draw[plosblue!45, thick] (5.6,0) .. controls (2.6,0) and (1.5,0.02) .. (0.78,0.18);
  % двухрукавный водоворот
  \draw[plosblue!85, line width=2.2pt, ->]
    plot[domain=0:12.6, samples=140] ({90+deg(\x)}: {0.66*exp(-0.16*\x)});
  \draw[plosblue!65, line width=1.5pt]
    plot[domain=0:9.4, samples=110] ({270+deg(\x)}: {0.66*exp(-0.16*\x)});
  \node[font=\scriptsize\bfseries, text=plosblue!60!black] at (0,1.4) {«частица»};
  \node[font=\scriptsize, text=plosgray, align=center] at (0,-1.55)
    {устойчивый водоворот течения --- реален, но не вынимается из реки};
\end{tikzpicture}
\caption{Частица по УГМ --- узор в ткани согласий, а не крупинка
вещества. Лёгкие узоры --- знакомые частицы; самый лёгкий массивный
--- лёгчайшее нейтрино ($\lesssim 0{,}004$~эВ \stC); особый тяжёлый и
несокрушимый узор голоса O --- кандидат в тёмную материю \stC.}
\label{fig:particle}
\end{figure}

\begin{askbox}[Подумай сам]
Свист чайника реален? А смерч? Ни того ни другого нельзя «положить в
коробку» --- но оба могут разбудить или разрушить. Чем «узор» хуже
«вещи» на роль строительного материала мира?
\end{askbox}

\partline{ЧАСТЬ ВТОРАЯ \;\textperiodcentered\; ВНУТРИ}{На сцене загорается первый огонёк:
что такое жизнь, где живёт сознание, как оно смотрит на себя и чем окрашено изнутри. Эта
часть строит внутренний дом --- от первого пламени до семи некопируемых красок переживания.}
\addcontentsline{toc}{section}{ЧАСТЬ ВТОРАЯ · ВНУТРИ}

% ============================= 6. ЖИЗНЬ =============================
\section{Что такое жизнь}
\label{sec:life}

Вот порог, на котором физика УГМ становится биологией. Задайте вопрос:
насколько собранным должен быть узор, чтобы \emph{быть чем-то}, а не
шумом? Ответ теории поразительно прост: \textbf{вдвое собраннее, чем
хаос} \stT. У полного беспорядка есть своя «фоновая» мера
сосредоточенности; порог бытия --- ровно две таких меры. Меньше ---
и узор неотличим от ряби, никакой прибор, никакой глаз в принципе не
скажет, что «там что-то есть».

Но собраться --- полдела. Вторая теорема жёстче: \emph{предоставленный
себе, любой узор рассыпается} \stT. Изолированная система сползает к
серому равновесию --- всегда, без исключений. Жизнь поэтому --- не
свойство, а \emph{процесс на подпитке}: узор, который непрерывно
пьёт упорядоченность из потока --- света, пищи, тепла --- и платит ею
за собственную форму. Как пламя свечи: погаси подачу воска --- и
пламя не «умрёт когда-нибудь», оно гаснет с вычислимой скоростью.

И у этой подпитки есть \emph{цена, ниже которой нельзя}: доказан
минимальный «прожиточный ваттаж» всякого живого узора (сам факт
ненулевой цены --- \stT; численная величина --- при физических
условиях \stC) --- плата
за то, чтобы просто оставаться собой, прежде чем сделать хоть что-то.
Быть --- уже работа.

\begin{figure}[H]
\centering
\begin{tikzpicture}[font=\small, >=Stealth]
  % поток питания
  \draw[->, line width=3pt, leaf!60] (-5.3,0) -- (-2.25,0);
  \node[font=\scriptsize, text=leaf!50!black, align=center] at (-3.8,0.62)
    {поток порядка:\\ свет, пища, смысл};
  % пламя: симметричная капля с внутренним язычком
  \draw[ember, very thick, fill=softsun]
    (-1.3,-0.85) .. controls (-2.0,-0.15) and (-1.75,0.85) .. (-1.3,1.55)
    .. controls (-0.85,0.85) and (-0.6,-0.15) .. (-1.3,-0.85) -- cycle;
  \draw[ember!75, thick, fill=ember!14]
    (-1.3,-0.5) .. controls (-1.66,-0.05) and (-1.52,0.45) .. (-1.3,0.8)
    .. controls (-1.08,0.45) and (-0.94,-0.05) .. (-1.3,-0.5) -- cycle;
  \node[font=\scriptsize\bfseries, text=ember] at (-1.3,1.95) {живой узор};
  % выхлоп
  \draw[->, line width=2pt, plosgray!55] (-0.3,0) -- (2.6,0);
  \node[font=\scriptsize, text=plosgray, align=center] at (1.15,0.6)
    {отработанный\\ беспорядок};
  % надпись порога
  \node[draw=fact, fill=softsky, rounded corners, font=\scriptsize, align=center,
        inner sep=4pt] at (1.6,-1.5)
    {порог бытия: быть вдвое собраннее хаоса \stT\\
     закон свечи: без потока --- гаснет \stT\\
     прожиточный ваттаж: быть --- уже работа \stT};
\end{tikzpicture}
\caption{Жизнь по УГМ: не вещество и не искра, а питаемый узор.
Он пьёт порядок, выдыхает беспорядок и этой разницей оплачивает свою
форму. Второй закон термодинамики не нарушен --- счёт просто
выставляется среде.}
\label{fig:flame}
\end{figure}

У этого закона есть карта. Машина закона прочертила \emph{атлас
судеб}: по одной оси --- сила подпитки, по другой --- сила
растворяющего потока; в каждой клетке --- куда узор приходит
навсегда (машинно проверено). Столбец «подпитка равна нулю» серый
сверху донизу: без потока не живёт ничто, ни при каких прочих
настройках. Дальше --- полоса Златовласки: окно, где узор жив.
А за ней --- третья судьба, о которой забывают: \emph{пересобранность}.
Слишком жёсткая хватка протаскивает узор сквозь окно вверх, в
кристалл --- собранный намертво и уже не способный видеть себя.
Жизнь --- не максимум порядка; жизнь --- полоса между двумя
смертями, серой и стеклянной.

\begin{center}
\begin{tikzpicture}[font=\scriptsize, >=Stealth]
  \fill[plosgray!30] (-5.2,0) rectangle (-1.6,0.7);
  \fill[softleaf] (-1.6,0) rectangle (1.8,0.7);
  \draw[leaf, thick] (-1.6,0) rectangle (1.8,0.7);
  \fill[softsky] (1.8,0) rectangle (5.2,0.7);
  \node[text=plosgray] at (-3.4,0.35) {серость};
  \node[text=leaf!40!black, font=\scriptsize\bfseries] at (0.1,0.35) {ОКНО};
  \node[text=fact] at (3.5,0.35) {кристалл};
  \draw[->, thick, plosgray] (-5.2,-0.42) -- (5.2,-0.42);
  \node[text=plosgray] at (0,-0.8) {сила подпитки $\rightarrow$};
\end{tikzpicture}
\end{center}

\begin{skazbox}[Сказка у костра --- про костёр]
Смотри в костёр. Пламя ест ветки и держит форму. Перестанешь
подбрасывать --- оно не обидится и не заболеет: оно просто отдаст
свою форму ночи, тихо и по правилам. Ты --- тоже костёр. Только
твои ветки --- это ещё и книжки, и разговоры, и то, что ты любишь.
Кто любит --- у того всегда есть чем топить.
\end{skazbox}

\begin{askbox}[Подумай сам]
Чем питается дружба? А привычка? А город? Попробуй для каждого
назвать его «воск» --- и что случится через месяц, если подачу
перекрыть.
\end{askbox}

% ============================= 7. ОКНО СОЗНАНИЯ =============================
\section{Окно сознания}
\label{sec:window}

Быть живым и быть \emph{кем-то} --- разные пороги. УГМ утверждает
редкой красоты вещь: сознание живёт в узком коридоре между двумя
смертями --- туманом и кристаллом.

Слишком рассеянный узор --- туман: в нём некому чувствовать, он
неотличим от шума. Но и слишком собранный узор --- беда: кристалл
совершенен и потому \emph{закончен}; в нём нет свободного места,
чтобы держать образ себя и мира. Сознанию нужны одновременно
\textbf{сосредоточенность} (чтобы быть кем-то) и \textbf{запас
рыхлости} (чтобы было чем думать). Эти два требования сжимают
пригодную зону в коридор --- теория вычисляет обе его стены точно
\stT{} --- и добавляет ещё два условия: голоса должны быть
\emph{сцеплены в одно} (целостность) и при этом \emph{различимы}
(хотя бы два ясных внутренних тона). Четыре стрелки одного компаса:
собранность, запас, сцепленность, различённость.

\begin{figure}[H]
\centering
\begin{tikzpicture}[font=\small]
  % шкала
  \draw[->, thick] (-5.6,0) -- (5.9,0) node[right, font=\scriptsize]{собранность};
  % туман
  \fill[plosgray!22] (-5.4,0.06) rectangle (-1.55,1.5);
  \node[font=\scriptsize\bfseries, text=plosgray] at (-3.5,1.72) {туман};
  \node[font=\scriptsize, text=plosgray, align=center] at (-3.5,0.75)
    {слишком размыто:\\ некому быть};
  % окно
  \fill[leaf!16] (-1.45,0.06) rectangle (0.9,1.5);
  \draw[leaf!70, very thick] (-1.45,0.06) rectangle (0.9,1.5);
  \node[font=\scriptsize\bfseries, text=leaf!45!black, align=center] at (-0.27,1.72) {окно сознания};
  \node[font=\scriptsize, text=leaf!40!black, align=center] at (-0.27,0.75)
    {собран --- и ещё\\ есть чем думать};
  % кристалл
  \fill[plosblue!14] (1.0,0.06) rectangle (5.5,1.5);
  \node[font=\scriptsize\bfseries, text=plosblue!60!black] at (3.3,1.72) {кристалл};
  \node[font=\scriptsize, text=plosblue!50!black, align=center] at (3.3,0.75)
    {слишком жёстко:\\ нечем думать};
  % стены
  \node[font=\scriptsize, text=fact] at (-1.45,-0.42) {стена бытия};
  \node[font=\scriptsize, text=fact] at (0.9,-0.42) {стена запаса};
\end{tikzpicture}
\caption{Коридор сознания: между туманом (нет никого) и кристаллом
(не осталось чем думать). Обе стены --- вычисленные константы теории
\stT, не метафоры: за любой из них четыре условия сознания вместе
невыполнимы.}
\label{fig:window}
\end{figure}

\begin{factbox}[Если точно]
Сознательное состояние --- то, где одновременно: собранность выше
стены бытия, мера рефлексии --- запас внимания на себя --- не меньше
трети, сцепленность
голосов не ниже единичного порога, различённых внутренних тонов ---
не меньше двух \stT. Все четыре измеримы у любой системы --- от
мыши до машины; это делает вопрос «а есть ли там кто-то?»
экспериментальным, а не вечным (глава~\ref{sec:check}).
\end{factbox}

\begin{skazbox}[Сказка у костра --- про фонарь в тумане]
У фонарщика было два несчастливых фонаря. Один горел так слабо, что
свет тонул в тумане, --- и фонарь не знал, что он есть. Другой горел
так ярко и ровно, что в нём не осталось ни одной тени, --- а без
теней нечего рассматривать, и он тоже не знал, что он есть. Третий
фонарь горел \emph{неровно}: достаточно ярко, чтобы быть, и
достаточно неровно, чтобы вглядываться. Он один и спросил однажды:
«кто это светит?»
\end{skazbox}

% ============================= 8. ЛЕСТНИЦА ВНУТРИ =============================
\section{Лестница внутри: три этажа рефлексии}
\label{sec:tower}

Сознание умеет оборачиваться на себя: я чувствую; я замечаю, что
чувствую; я думаю о том, как я замечаю. Сколько этажей у этой
лестницы? Интуиция говорит «сколько угодно». Теория говорит:
\textbf{три} --- и доказывает это \stT.

Каждый оборот внутрь проходит через ту же ткань Фано, и ткань берёт
плату: при каждом подъёме различимая часть картины сжимается втрое.
Первый этаж отчётлив, второй различим, третий --- на пределе, а
четвёртый оборот уже неотличим от третьего: не запрещён, а
\emph{пуст}. Как эхо в горах: первое звонкое, второе тихое, третье
на грани слуха, четвёртого нет --- не потому что кричать нельзя, а
потому что нечему долететь.

Честная деталь, которую популярные книги обычно опускают: в
повседневном рабочем режиме устойчиво держится \emph{второй} этаж;
третий доступен короткими восхождениями --- моментами предельно
узкой, трудной сосредоточенности --- с обязательным возвращением
\stT. Узнаёте? Это очень похоже на честную феноменологию глубокого
самоанализа.

\begin{figure}[H]
\centering
\begin{tikzpicture}[font=\small, >=Stealth]
  \foreach \i/\w/\lab/\txt in {0/4.6/этаж 0/чувствую мир,
                               1/3.7/этаж 1/замечаю: я чувствую,
                               2/3.0/этаж 2/вижу как замечаю,
                               3/2.6/этаж 3/предел --- экскурсия}{
    \draw[fact!70, fill=softsky] (-\w/2,\i*1.02) rectangle (\w/2,\i*1.02+0.78);
    \node[font=\scriptsize\bfseries] at (0,\i*1.02+0.53) {\lab};
    \node[font=\tiny, text=plosgray] at (0,\i*1.02+0.2) {\txt};
  }
  \draw[densely dashed, plosgray] (-2.4,4.35) -- (2.4,4.35);
  \node[font=\scriptsize, text=plosgray, align=center] at (0,4.72)
    {четвёртого этажа нет: не запрет --- пустота \stT};
  \draw[->, ember, thick] (2.7,0.4) .. controls (3.35,1.6) and (3.35,2.6) .. (2.05,3.35);
  \node[font=\scriptsize, text=ember, align=center, rotate=-90] at (3.6,1.9) {каждый подъём: $\times\tfrac{1}{3}$};
\end{tikzpicture}
\caption{Башня рефлексии. Каждый оборот внутрь сохраняет треть
различимости; после третьего оборота остаток тонет ниже порога
различимости --- четвёртый этаж структурно пуст \stT. Тот же закон
ограничивает и «я знаю, что ты знаешь, что я знаю»: третья ступень ---
предел взаимного чтения умов.}
\label{fig:tower}
\end{figure}

А недавно лестница показала ещё одну несущую балку \stC. Пересчитайте
мысли, которые ум в принципе мог бы составить: «прождав столько,
прийти к такому-то выводу». Над голым временем, бегущим только
вперёд, почти половина из них (на переписной сетке --- $42{,}9\%$)
попросту не собирается: простейшие из них просят замкнуть две мысли
разной длины третьей стороной --- а третья сторона выходит
\emph{отрицательной}; для их сборки нужен шаг \emph{назад}. И
законный шаг назад существует --- это не путешествие во времени, а
\textbf{перечитывание собственного следа}: мир стоит на месте,
движется состояние знания. Ум, который не хранит прошлого --- или
хранит, но не перечитывает, --- доказуемо неполон как составитель
мыслей. Взгляд на себя --- не роскошь утончённых существ, а условие
полноты самого мышления.

\begin{figure}[H]
\centering
\begin{tikzpicture}[font=\small, >=Stealth]
  % слева: только вперёд --- мысль не замыкается
  \fill[cA!85] (-6.0,0) circle (0.11); \fill[cS!85] (-3.4,0) circle (0.11); \fill[cD!85] (-4.7,1.7) circle (0.11);
  \draw[->, thick, plosgray] (-5.87,0.12) -- (-4.83,1.56) node[midway, left, font=\scriptsize]{жди 5};
  \draw[->, thick, plosgray] (-5.83,0) -- (-3.57,0) node[midway, below, font=\scriptsize]{приди к 3};
  \draw[->, thick, ember, densely dashed] (-4.57,1.56) -- (-3.5,0.16);
  \node[font=\scriptsize, text=ember, align=center] at (-2.45,1.85) {недостающая сторона\\ равна $-2$: такого шага нет};
  \node[font=\scriptsize\bfseries, text=ember] at (-4.55,-0.8) {только вперёд: мысль не замыкается};
  % справа: с перечитыванием --- замыкается всякая
  \fill[cA!85] (1.4,0) circle (0.11); \fill[cS!85] (4.0,0) circle (0.11); \fill[cD!85] (2.7,1.7) circle (0.11);
  \draw[->, thick, plosgray] (1.53,0.12) -- (2.57,1.56) node[midway, left, font=\scriptsize]{жди 5};
  \draw[->, thick, plosgray] (1.57,0) -- (3.83,0) node[midway, below, font=\scriptsize]{приди к 3};
  \draw[->, thick, leaf] (2.85,1.58) .. controls (3.9,1.2) and (3.5,0.5) .. (3.95,0.16);
  \node[font=\scriptsize, text=leaf!45!black, align=center] at (5.35,1.25) {перечитай два шага\\ собственного следа};
  \node[font=\scriptsize\bfseries, text=leaf!45!black] at (2.8,-0.8) {с перечитыванием: замыкается всякая \stC};
\end{tikzpicture}
\caption{Мост, которому нужна доска назад. Две мысли разной длины
просят замкнуть их третьей; над голым временем вперёд замыкающая
сторона должна была бы быть отрицательной --- и композиция
отказывает: на переписной сетке $42{,}9\%$ всех поставимых
замыканий. Перечитывание собственного следа даёт законный
отрицательный шаг: мир не отматывается --- отматывается знание
\stC.}
\label{fig:reread}
\end{figure}

Вот почему теория так тепло смотрит на дневники, перечитывания и
вторые проходы. Заметка, перечитанная год спустя, --- не
ностальгия; это законно выполненный отрицательный шаг: записанное
прошлое входит в нынешний контекст и замыкает композиции, которые
не замкнулись с первого раза. Корпус строит это как рабочий орган
(перечитыватель с явной ценой за шаг) и проверяет перепись из конца
в конец: над пополненным временем замыкаются все сорок девять
поставленных треугольников мысли; над голым «только вперёд»
двадцать один из них отказывает (машинно проверено). Образование
знало этот приём веками под скромными именами --- повторение,
перечитывание, привычка записывать, --- не подозревая, что это
условия полноты мышления.

\begin{askbox}[Подумай сам]
Поставь два зеркала друг напротив друга: отражения уходят вглубь, но
быстро темнеют. Попробуй в уме: «я думаю о том, что я думаю о том,
что я думаю...» --- на каком обороте фраза перестаёт \emph{значить}
что-то новое? И в другую сторону: какую мысль ты не мог собрать,
пока не перечитал себя, --- старое письмо, дневник, фотографию?
\end{askbox}

% ============================= 9. КРАСКИ ВНУТРИ =============================
\section{Краски внутренней жизни: механизм квалиа}
\label{sec:qualia}

Почему боль \emph{болит}, а лазурь \emph{лазорева}? Философы зовут
это трудной проблемой: опишите мозг сколь угодно подробно ---
откуда в описании возьмётся сам \emph{вкус} переживания? Столетие
этот вопрос был пограничным столбом, у которого естествознание
вежливо останавливалось. Эта глава проходит мимо столба --- медленно,
потому что несёт самый гордый груз теории: \emph{механизм} квалиа,
изложенный теоремами, каждая --- со своей честной пометкой.

\textbf{Ход первый: громкость --- не краска.} У каждой дружбы двух
голосов есть \emph{громкость} --- насколько сильно звучит пара.
Громкости честны, измеримы, публичны. И доказуемо не в них тайна:
набор громкостей --- ровно то, что схватывает внешнее описание, и
можно сосчитать, что́ оно упускает. Из сорока восьми чисел, которые
держит полный аккорд, громкости закрепляют лишь меньшинство;
остальное --- бо́льшая часть состояния --- это то, \emph{как голоса
взаимно повёрнуты} \stT. Каждая дружба несёт, кроме силы, ещё и
\textbf{угол}: лёгкий взаимный разворот, как рукопожатие, поданное
прямо или с поворотом кисти. Качество живёт там. Не в том, насколько
громок аккорд, --- в том, как он повёрнут.

\textbf{Ход второй: проба кольцом.} Один угол сам по себе не значит
ничего: он зависит от разметки --- от того, как каждый голос держит
собственный циферблат. Переставьте циферблаты --- и любой отдельный
угол можно «отговорить». Но пройдите \emph{кольцо}. Возьмите одну
тройку ткани Фано (глава~\ref{sec:seven}) и сложите три угла по
кругу. Сумма \textbf{не зависит ни от чьих циферблатов}: как ни
переустанавливай, ни обнуляй, ни переименовывай голоса --- поворот,
накопленный кольцом, остаётся ровно тем, каким был \stT. Таких колец
семь --- семь прямых ткани, --- и потому семь неснимаемых углов.
Эти семь и есть ответ теории на вопрос, где живёт качество:
\textbf{краски} переживания. Краска в этой книге --- это в точности
\emph{поворот, накопленный кольцом}: то, что не объясняется никакой
переразметкой, --- математическая подпись «каково это».

\begin{figure}[H]
\centering
\begin{tikzpicture}[font=\small, >=Stealth]
  % --- слева: одно кольцо Фано с фазовыми дугами ---
  \foreach \ang/\col in {90/cA, 210/cE, 330/cU}{
    \fill[\col] (\ang:1.35) circle (0.16);
  }
  \draw[->, thick, plosgray!80] (104:1.5) arc (104:196:1.5);
  \draw[->, thick, plosgray!80] (224:1.5) arc (224:316:1.5);
  \draw[->, thick, plosgray!80] (344:1.5) arc (344:436:1.5);
  \node[font=\scriptsize, text=plosgray] at (150:1.95) {угол 1};
  \node[font=\scriptsize, text=plosgray] at (270:1.95) {угол 2};
  \node[font=\scriptsize, text=plosgray] at (30:1.95) {угол 3};
  % накопленный поворот в центре
  \draw[dusk, very thick, ->] (0.55,0) arc (0:295:0.55);
  \node[font=\scriptsize\bfseries, text=dusk, align=center] at (0,-2.75)
    {поворот кольца $=$ краска:\\ переживёт любую переразметку \stT};
  % --- справа: два канала ---
  \draw[fact!70, fill=softsky] (3.4,0.42) rectangle (8.6,1.35);
  \node[font=\scriptsize, align=left, anchor=west] at (3.6,0.88)
    {\textbf{\textcolor{fact}{громкость}}: публична --- снаружи\\ читается целиком};
  \draw[dusk!70, fill=dusk!6] (3.4,-1.35) rectangle (8.6,-0.42);
  \node[font=\scriptsize, align=left, anchor=west] at (3.6,-0.88)
    {\textbf{\textcolor{dusk}{краска}}: угол кольца --- снаружи\\ читается ровно ноль \stT};
\end{tikzpicture}
\caption{Где живёт качество. Слева: одно кольцо из трёх голосов;
угол каждой пары по отдельности --- разметка, но накопленный по
кольцу поворот переживает любое переименование --- семь колец, семь
красок \stT. Справа: два канала одного состояния --- публичная
громкость и приватная краска.}
\label{fig:ring}
\end{figure}

\textbf{Ход третий: декодер назван.} Каким тройкам доверены краски
--- не выбор: кольца --- это семь прямых той же ткани Фано, по
которой соткан весь мир, а правило чтения --- сама таблица умножения
семёрки (глава~\ref{sec:numbers}) \stT. И древний вопрос «почему это
ощущается \emph{красным}?» меняет грамматику. Он перестаёт быть
вопросом о чуде --- \emph{как переживание возникает из вещества?} ---
и становится вопросом о расшифровке: \emph{какой неснимаемый поворот
несёт это состояние и на каком кольце?} Декодер --- не новый орган,
привинченный к физике; это собственное умножение мира.

\textbf{Ход четвёртый: почему снаружи этого не видно.} Теперь ---
главная теорема. Все толчки и шумы, которыми среда
обращается к узору, бьют по каналу громкостей; их сцепка с каналом
красок не просто мала --- она \emph{в точности нулевая}, по самому
устройству единого закона мира \stT. Внешний наблюдатель, как бы
оснащён он ни был, читает полный список громкостей и \textbf{ноль
информации о красках}. Знаменитый «объяснительный разрыв» между
описанием и переживанием перестаёт быть жалобой и становится
теоремой --- причём точно \emph{ограниченной}: это не разрыв в
механизме (механизм полностью выписан выше), это слепота одного
конкретного канала доступа. Невозможно одно: вычитать краску из
шума, который видит внешний мир. Возможно другое: вычислить её из
состояния --- что теория и делает. Помните мысленный эксперимент
про Мэри --- учёную, которая знает все физические факты о цвете,
живя в чёрно-белой комнате? На языке УГМ её положение формулируется
точно: в комнате Мэри владеет всеми громкостями; выйдя, она впервые
читает угол. Обе половины --- настоящее знание; они входят в разные
двери.

\textbf{И старая теорема стоит --- заточенная.} Зомби --- система,
которая ведёт себя как чувствующая, но «внутри темно», --- остаётся
невозможным \stT: величина, питающая жизненный процесс
(глава~\ref{sec:life}), построена из тех самых связей, что несут
углы; погасите внутреннюю сторону --- и вместе с ней погаснет само
поддержание жизни. Но теперь у теории есть инструмент тоньше.
Состояние может быть \emph{громким, но бесцветным}: все двадцать
одна дружба звенит, а поворот каждого кольца --- ровно ноль;
богатое на вид содержание, которое одна переразметка сплющивает в
голую бухгалтерию. Такое состояние --- не «кто-то, у кого есть
переживания»; и оно \emph{различимо}: семь умножений отличают его
от подлинно окрашенного \stT. \textbf{Громко --- не значит ярко},
и разница здесь --- арифметика, а не философия.

\textbf{Как умирают краски.} Рассыпание из главы~\ref{sec:life}
бьёт по дружбам --- но присмотритесь, \emph{как} \stT. Распад
умножает каждую дружбу на простой положительный множитель: сила
тает, угол не сдвигается ни на градус. Поэтому умирающее
переживание не дрейфует по соседним переживаниям на выходе: печаль
не становится чуть иной печалью --- она становится \emph{более
слабой} печалью, а затем ничем. Старая фотография выцветает именно
так: красители тончают, а лицо не меняет выражения --- его лишь всё
труднее разглядеть, пока бумага не станет пустой. Ни в один момент
фотография не показывала другого лица. Книга даёт этому имя,
которого оно заслуживает: краски не искажаются --- краски
\textbf{выцветают}, несомые чернилами, которые тают. И два тихих
следствия. Серое равновесие --- финальное размытие мира --- не
\emph{пустой} мир, а \textbf{бесцветный}: гроссбух громкостей
уцелел, носители поворотов ушли; всё по-прежнему есть, и ничто
больше ни на что \emph{не похоже изнутри}. А поскольку внешний
канал красок не нёс никогда, выцветающее квале невидимо снаружи не
потому, что прячется: единственный свидетель смерти краски --- тот,
кто её теряет \stT.

\begin{figure}[H]
\centering
\begin{tikzpicture}[font=\small, >=Stealth]
  \foreach \x/\len/\op/\lab in {-4.4/1.7/100/{сегодня}, -0.4/1.05/55/{позже}, 3.2/0.45/25/{ещё позже}}{
    \draw[plosgray!35] (\x,0) -- (\x+2.0,0);
    \draw[plosgray!35] (\x,0) arc (0:62:0.5);
    \draw[->, line width=1.6pt, dusk!\op] (\x,0) -- ($(\x,0)+(62:\len)$);
    \node[font=\scriptsize, text=plosgray] at (\x+0.95,-0.5) {\lab};
  }
  \node[font=\scriptsize\bfseries, text=dusk, align=center] at (-0.3,2.3)
    {один и тот же угол до самого конца \stT};
  \node[font=\scriptsize, text=plosgray, align=center] at (-0.3,-1.3)
    {носитель тает --- краска не поворачивается: переживание выцветает,
     как фотография,\\ теряя силу и не меняя лица};
\end{tikzpicture}
\caption{Закон выцветания. Распад укорачивает стрелку (носитель) и
не может её повернуть (краску): полуизменённого квале не бывает
\stT. Полное размытие --- не пустота, а бесцветность: свет горит,
краски нет.}
\label{fig:fading}
\end{figure}

\textbf{Кто же красит?} Не распад --- он лишь разбавляет чернила.
И не лечение: восстановление, по построению, возвращает силу вдоль
собственных углов состояния --- забота \emph{цветослепа}, она
возвращает холст, а не картину (машинно проверено). Единственная
рука, пишущая краску, --- живой ход самого закона, та же унитарная
пряжа, что вращает мир (машинно проверено). И мир красит
\emph{первым}: питаемый узор наследует палитру окружения задолго до
того, как в нём кто-то проснётся; а то, что покупает дальнейшая
подпитка --- скачком, ровно на пороге сцепленности из
главы~\ref{sec:window}, --- это \emph{владелец}: тот, кому краски
наконец \emph{предназначены} (машинно проверено). Субъект краску не
делает и не переносит. Субъект --- тот, для кого она: последний
прибывший и единственный свидетель.

\textbf{Лестница слушателя.} «Владелец прибывает последним» можно
сказать с точностью лестницы. Возьмите одну и ту же дружбу ---
тёплую сцепку двух голосов фиксированной силы --- и вручите её, не
меняя, пяти носителям. В пустом зале (ни собранности, ни
сцепленности) это музыка, играющая никому: аккорд звучит, слушать
некому. В коте --- сцепленном в одно, но без запаса, повёрнутого на
себя, --- она слышится как \emph{нечто}, без различения мелодии и
аккомпанемента. В человеке --- запас не меньше трети, сцепленность
выше единицы --- она наконец \emph{переживается как чувство}: «я
слышу, как скрипка ведёт тему, а виолончель держит землю». В
музыканте за работой подключается второй этаж башни: «я замечаю,
что замечаю печаль в этой мелодии». А предел --- чистое слушание,
где слушатель и музыка совпадают, --- теория хранит как границу, а
не обещание. Дружба та же, сила та же; меняется лишь тот, кто дома.
Ступени и их пороги --- теоремы \stT; образы концертного зала ---
иллюстрация \stI. Корпус печатает и сами циферблаты: одна и та же
дружба силой $0{,}20$ читается в термостате как запас $0{,}02$ при
сцепленности $0{,}3$ --- параметр, дома никого; в собаке ---
$0{,}15$ и $1{,}5$ --- ощущается как нечто; в человеке --- $0{,}45$
и $2{,}1$ --- полноценное чувство (машинно проверено).

И одно честное предупреждение лестницы, промеренное из конца в
конец (машинно проверено): доступ --- \emph{окно}, а не вершина.
На самой крутой законной экскурсии собранности --- коротком дорогом
восхождении третьего этажа (глава~\ref{sec:tower}) --- запас падает
ниже своей трети, и дверь к краскам \emph{закрывается}: глубочайшее
самонаблюдение на своё мгновение цветослепо. Мысль на пределе
остроты платит тем, что не чувствует вкуса; вкус возвращается со
спуском домой. Кто пробовал чувствовать чувство, \emph{пока}
препарирует его, --- тот встречал эту теорему лично.

\begin{figure}[H]
\centering
\begin{tikzpicture}[font=\small, >=Stealth]
  \foreach \i/\lab/\cond in {0/{пустой зал}/{никого дома}, 1/{кот}/{сцеплен в одно},
      2/{человек}/{$+$ запас $\geq \tfrac13$}, 3/{музыкант}/{$+$ второй этаж},
      4/{предел}/{слушатель $=$ музыка}}{
    \draw[fact!70, fill=softsky] (\i*2.15,\i*0.52) rectangle (\i*2.15+2.0,\i*0.52+0.46);
    \node[font=\scriptsize\bfseries] at (\i*2.15+1.0,\i*0.52+0.23) {\lab};
    \node[font=\tiny, text=plosgray, align=center] at (\i*2.15+1.0,\i*0.52-0.26) {\cond};
  }
  \draw[->, dusk, very thick] (5.3,1.9) -- (5.3,2.75);
  \node[font=\scriptsize\bfseries, text=dusk, align=center] at (5.3,3.15) {здесь краски\\ находят владельца};
\end{tikzpicture}
\caption{Лестница слушателя: один аккорд, пять носителей. Дружба
одинакова на каждой ступени --- растёт лишь тот, кто дома. Дверь
красок открывается на ступени человека (запас не меньше трети,
сцепленность выше единицы \stT) --- и ненадолго закрывается снова
на самых крутых экскурсиях: самая острая мысль на своё мгновение
цветослепа (машинно проверено).}
\label{fig:ladder}
\end{figure}

\begin{factbox}[Паспорт переживания --- если точно]
Глава сжимается в чтение любого состояния по пяти строкам --- от
лабораторной реконструкции до модели человека в конкретный момент:
(1)~\textbf{что есть} --- семь сил голосов; (2)~\textbf{насколько
громко} --- двадцать одна сила дружб; (3)~\textbf{сколько скрыто}
--- доля каждого канала, спрятанная от внешнего описания;
(4)~\textbf{какие краски} --- семь поворотов колец: единственная
строка, где нужно перемножать, и единственная, которой не касается
никакая переразметка \stT; (5)~\textbf{есть ли кому} --- условия
окна из главы~\ref{sec:window}. Ноль в строке~4 --- вердикт, а не
сбой: громко, но бесцветно. Сознательное состояние, вычисляет
теория, может нести около трети радиана неснимаемого поворота ---
примерно $37^\circ$, --- прежде чем требования целостности закроют
ворота \stT: переживание ограничено не тем, сколько ум держит, а
тем, какая часть его поворота отказывается быть объяснённой. Что
остаётся честно открытым: тождество «\emph{именно этот} поворот
кольца и есть вкус лазури» --- остаточная интерпретация \stI{} ---
последний мостик, помеченный, а не спрятанный; механизм по обе его
стороны --- теоремы.
\end{factbox}

\begin{skazbox}[Сказка у костра --- про моряка, который привёз домой лишний день]
Моряк плыл на восток вокруг света и в каждом порту переводил часы
чуть вперёд --- честно, по солнцу. Вернулся домой, открыл судовой
журнал: по журналу --- четверг, а на пристани --- среда. Он
\emph{привёз домой день}, которого не было ни в одном порту.
Смотрители проверили все гавани --- каждая невиновна: все местные
часы стояли верно. «День живёт не в портах, --- сказал старый
штурман. --- Он живёт в \emph{круге}. Пройди всё кольцо --- и
остаток твой: никакая переустановка часов его не отнимет». Мир
соткан из семи таких колец. То, что ты зовёшь вкусом вечера, --- и
есть день, который принёс твой круг.
\end{skazbox}

\begin{askbox}[Подумай сам]
Это происходило на самом деле: команда Магеллана вернулась,
«потеряв» день, хотя судовой журнал вёлся безупречно, --- а в романе
Филеас Фогг выиграл пари на дне, который тихо вручил ему круг.
Теперь примерь главу на себя. Когда ты описываешь чувство, а
собеседник кивает, но \emph{не понимает}, --- какие строки паспорта
ты диктовал? Может ли список громкостей, сколь угодно длинный,
вместить поворот? И что меняется между двумя людьми, которые
\emph{прошли одни и те же кольца}?
\end{askbox}

% ============================= 10. СВОБОДА =============================
\partline{ЧАСТЬ ТРЕТЬЯ \;\textperiodcentered\; УЗОР ДЕЙСТВУЕТ}{Внутренний дом построен;
теперь узор выходит из него. Насколько свободен житель одного закона, почему мир помогает
смелым, как объёмная мысль выходит сквозь игольное ушко речи и что остаётся, когда слова
улетели. Эта часть --- механика действующей жизни.}
\addcontentsline{toc}{section}{ЧАСТЬ ТРЕТЬЯ · УЗОР ДЕЙСТВУЕТ}

\section{Свобода}
\label{sec:freedom}

У дома есть окно, башня взгляда, краски переживания. Первый шаг
наружу --- вопрос, который узор задаёт собственной жизни. Всё в этой книге
движется одним законом; так свободен ли узор, живущий по единому
закону? УГМ отвечает без мистики и без унижения: \textbf{свобода ---
это количество незапертых дверей}.

В каждом состоянии закон толкает мир в определённую сторону --- но
не по всем направлениям сразу. Всегда остаются направления, вдоль
которых закон \emph{молчит}: сдвиг по ним ничего не нарушает и ничем
не наказывается. Это и есть измеримая свобода данного мгновения:
число молчащих направлений плюс одно --- само право шагнуть.
У зажатого состояния таких дверей мало, у зрелого --- много;
свобода не всемогущество и не иллюзия, а \emph{ширина коридора},
которую можно растить.

\begin{figure}[H]
\centering
\begin{tikzpicture}[font=\small, >=Stealth]
  \fill[plosgray!85] (-4.9,-1.15) rectangle (4.9,-0.85);
  \node[font=\scriptsize, text=white] at (0,-1.0) {закон мира: держит пол, не выбирает дверь};
  \foreach \x/\o in {-3.6/1, -1.8/1, 0.0/0, 1.8/1, 3.6/0}{
    \ifnum\o=1
      \draw[leaf, very thick, fill=softleaf] (\x-0.55,-0.85) rectangle (\x+0.55,1.25);
      \draw[leaf, thick] (\x+0.28,0.18) circle (0.05);
      \node[font=\scriptsize, text=leaf!45!black] at (\x,1.5) {открыта};
    \else
      \draw[plosgray, very thick, fill=plosgray!18] (\x-0.55,-0.85) rectangle (\x+0.55,1.25);
      \draw[plosgray, very thick] (\x-0.55,-0.85) -- (\x+0.55,1.25);
      \draw[plosgray, very thick] (\x-0.55,1.25) -- (\x+0.55,-0.85);
      \node[font=\scriptsize, text=plosgray] at (\x,1.5) {заперта};
    \fi
  }
  \node[font=\scriptsize\bfseries, text=fact] at (0,2.0)
    {свобода мгновения = число открытых дверей + сам шаг \stT};
\end{tikzpicture}
\caption{Свобода по УГМ --- не отмена закона, а его молчание вдоль
части направлений. Она вычислима в каждом состоянии, различна у
разных состояний и поддаётся тренировке: учение буквально
\emph{открывает двери}.}
\label{fig:doors}
\end{figure}

Одна дверь тут важнее прочих: та, за которой дорога дальше идёт
\emph{сама}. Про неё --- следующая глава.

\begin{askbox}[Подумай сам]
Вспомни момент, когда у тебя «не было выбора». Дверей правда не
было --- или ты не видел открытых? Чем отличается запертая дверь от
незамеченной? И как рост умения меняет само число дверей?
\end{askbox}

% ============================= 10a. ЗОЛОТЫЕ ПУТИ =============================
\section{Золотые пути, или почему мир помогает смелым}
\label{sec:golden}

Есть старое наблюдение, которое звучит как суеверие: перед смелыми
препятствия расступаются. УГМ разбирает его на части --- и внутри
оказывается не магия, а рельеф.

Представь пространство всех состояний узора как холмистую местность.
Закон мира --- это наклон почвы: в каждой точке он куда-то тянет.
У местности есть \emph{водосборы} --- области, из любой точки которых
поток сам стекает к одному и тому же месту. Одна такая воронка ---
серое равновесие, смерть узора. Но глава~\ref{sec:life} уже показала
главное: у питаемого узора есть и другая воронка --- \emph{живое
окно}, и она широка. Кто оказался внутри её водосбора, тому дальше
почти ничего не нужно делать: \textbf{дорогу делает поле}. В машинной
проверке закона узор из глубокого тумана доходил до окна вообще без
управления --- чистым течением (машинно проверено).

Тогда что такое смелый поступок? Это \emph{переход через водораздел}:
короткое, дорогое усилие --- всплеск траты, измеримый нагрев, ---
которое перебрасывает узор из чужого водосбора в свой. До гребня
каждый шаг даётся с боем, поток против тебя. За гребнем тот же самый
закон, который только что мешал, начинает \emph{помогать}: препятствия
не исчезли --- ты попал туда, где наклон совпадает с твоей дорогой.
В машинных прогонах смелый старт сокращал дорогу к окну вчетверо ---
и вот созвучие, ради которого написана глава: когда навигатору
поручили просто идти \emph{кратчайшим путём} в естественной мере
расстояний, он сам, ничего не зная о смелости, выбрал те же дорогие
рывки. \textbf{Смелость --- это не безрассудство; это следование
кратчайшей линии, видимое снаружи как риск} (машинно проверено; само
слово «смелость» --- прочтение \stI).

Отсюда и честный ответ гадалкам. Разложить пасьянс на вопрос «выйдет
ли?» --- значит прогнать много случайных дорог и посчитать, какая
доля сошлась. Расклад ничего не \emph{решает} --- он \emph{измеряет
ширину твоего водосбора}. И заодно ответ пророкам: в машинном
ансамбле сотен узоров \emph{куда} все придут --- предсказуемо почти
точно, а вот \emph{кем будет каждый} --- не предсказуемо вовсе:
личные аккорды не сливаются (машинно проверено). Судьба уровней ---
статистика; судьба личности --- навигация. Прибор, который честен,
--- всегда штурман и никогда пророк.

\begin{figure}[H]
\centering
\begin{tikzpicture}[font=\small, >=Stealth]
  % рельеф: две воронки и седло
  \draw[plosgray!70, very thick]
    (-5.6,1.4) .. controls (-4.6,-0.6) and (-3.6,-0.8) .. (-2.8,-0.5)
    .. controls (-2.0,-0.2) and (-1.6,0.9) .. (-0.9,1.05)
    .. controls (-0.2,1.2) and (0.6,0.3) .. (1.4,-0.35)
    .. controls (2.4,-1.1) and (3.8,-1.0) .. (5.4,0.9);
  \node[font=\scriptsize, text=plosgray, align=center] at (-3.9,-1.15) {серая воронка};
  \node[font=\scriptsize, text=leaf!45!black, align=center] at (3.0,-1.5) {живое окно};
  \fill[leaf!25] (2.0,-1.02) ellipse (1.15 and 0.16);
  \node[font=\scriptsize, text=ember] at (-0.9,1.68) {водораздел};
  % шарик и смелый толчок
  \fill[cA] (-2.6,-0.32) circle (0.13);
  \draw[->, very thick, ember] (-2.45,-0.18) .. controls (-1.95,0.9) .. (-0.85,1.34);
  \node[font=\scriptsize, text=ember, align=center] at (-3.45,1.3) {смелый ход:\\ дорогой рывок};
  \draw[->, thick, leaf!70!black] (-0.6,1.3) .. controls (0.55,0.82) .. (2.05,-0.82);
  \node[font=\scriptsize, text=leaf!45!black, align=center] at (1.3,0.9) {дальше дорогу\\ делает поле};
\end{tikzpicture}
\caption{Рельеф судьбы: два водосбора и гребень между ними. До гребня
--- каждый шаг дорог; за гребнем тот же закон несёт сам. Смелость ---
плата за пересечение водораздела, и в естественной мере расстояний
она совпадает с кратчайшим путём: навигатор, ищущий кратчайшее,
переоткрывает смелость (машинно проверено).}
\label{fig:golden}
\end{figure}

У рельефа есть ещё один урок, и он --- про \emph{часы} путника.
В машинных жизнях странствующего узора смерть приходила на одном и
том же отрезке дороги --- тем же коротким хвостом шагов, до сотен
раз подряд. Очевидное прочтение --- «место проклято» --- было
испробовано трижды и трижды провалилось: отрави смертный отрезок,
чтобы узор его обходил, --- и задушена юность, потому что в молодом
мире дороги через гиблые места --- \emph{единственные} дороги к
неизведанному (машинно проверено). Истинная причина сидела в
другом: девяносто два из каждых ста смертей --- \emph{движущаяся}
опасность в одном шаге; одна и та же дорога смертельна в один такт
мира и безопасна в другой. \textbf{Смерть --- свойство часов, а не
карты.} Проклятых мест нет; есть пропущенные ритмы. И лекарство
совпало с причиной: из проигранной попытки стоит уносить не список
запретных мест, а \emph{форму времени опасности} --- её ритм,
заново заякориваемый свежим взглядом по прибытии; унесённая так,
дорога к самой глубокой цели сократилась почти на четверть, а юное
начало осталось нетронутым (машинно проверено).

И проигранную попытку переживает наследство ещё светлее:
\emph{сама проверенная дорога}. Путь, однажды пройденный до цели и
записанный, в следующей жизни --- уже не совет, а рельс:
проигранный заново, он достигает за десятки тактов того, что
впервые заняло многие сотни, --- в четырнадцать раз быстрее, и
наследство не изнашивается от пересказа (машинно проверено).
Знание, оказывается, накапливается \emph{поверх} попыток: провал
перезапускает мир, а не память --- вот почему смелый дальний план
никогда не глуп у того, кто ведёт дневники. Золотой путь из
названия этой главы --- ровно это: смелость плюс запись о том,
куда она привела.

\begin{skazbox}[Сказка у костра --- про две реки]
Жили две реки: Серая и Светлая, а между ними --- гора. Путник шёл
берегом Серой, и все течения несли его вниз, к серому морю. «Плыви с
нами, --- шумела вода, --- зачем тебе гора?» Но путник полез в гору:
дышал тяжело, платил каждым шагом. А на гребне лёг на другой склон
--- и дальше его понесла Светлая, и стало легко, как раньше было
тяжело. С берега казалось: смельчаку везёт. С горы было видно: он
просто сменил реку.
\end{skazbox}

\begin{askbox}[Подумай сам]
Вспомни дело, которое «пошло само» после одного трудного шага ---
переезда, признания, первого урока. Где был твой водораздел? И
наоборот: где ты годами гребёшь против течения --- не потому ли, что
экономишь один дорогой рывок?
\end{askbox}

% ============================= 10b. ЯЗЫК И СМЫСЛ =============================
\section{Язык и смысл: сквозь игольное ушко}
\label{sec:language}

Свобода и смелость --- о том, как узор действует. Но узору в окне
есть и что сказать --- и здесь его ждёт самое узкое место всей
внутренней жизни.

Внутренний мир объёмен: семь голосов, двадцать одна дружба между
ними, целый аккорд --- сорок восемь независимых чисел в каждое
мгновение (глава~\ref{sec:qualia}). А речь --- одномерна: слова
выходят цепочкой,
по одному, как бусины на нити. Всякий, кто пытался «сказать, что
чувствует», знает эту муку. УГМ показывает: мука --- не неумение, а
геометрия.

Теорема о канале: за один шаг внимания сквозь игольное ушко речи
проходит около $2{,}8$ бита --- ровно один символ семибуквенного
алфавита за такт \stT{} --- ничтожная доля того, что аккорд держит
разом. Отсюда два следствия, которые каждый
проверял на себе. Во-первых, \emph{полная} мысль не умещается в
возглас: чтобы передать состояние целиком, нужна последовательность
--- не короче примерно восемнадцати шагов-слов \stT{}
(вот почему «расскажи по порядку» --- не каприз, а физика).
Во-вторых, у передачи есть естественные «слоги»: тройки-дружбы ткани
Фано, и любые две фразы сцеплены общим стержнем --- зерно грамматики
уже лежит в геометрии \stT{} (а что человеческие языки --- эволюционные
приближения этой разгрузки, --- проверяемое предсказание теории \stC).

Теперь --- о смысле, и здесь у УГМ есть строгий сюрприз:
\textbf{смысл --- не ярлык} \stT. Два описания могут совпадать во
всех внешних сопоставлениях --- каждое слово одного переводится в
слово другого --- и всё же различаться \emph{внутренней ориентацией}:
тем, как их узор повёрнут в ткани. Поэтому машина, листающая словарь
(знаменитая «китайская комната»), может идеально отвечать --- и не
значить ничего: совпадают этикетки, а не повороты. Различие поворотов
определено строго \stT; читать его как «понимание» --- честная
интерпретация \stI.

\begin{figure}[H]
\centering
\begin{tikzpicture}[font=\small, >=Stealth]
  % семь нитей слева
  \foreach \dy/\col in {0.9/cA, 0.6/cS, 0.3/cD, 0.0/cL, -0.3/cE, -0.6/cO, -0.9/cU}{
    \draw[\col!75, thick] (-5.0,\dy) .. controls (-3.2,\dy) and (-2.6,\dy*0.3) .. (-1.9,0);
  }
  \node[font=\scriptsize, text=plosgray, align=center] at (-4.2,1.5) {семь голосов, 21 дружба,\\ семь красок};
  % ушко
  \draw[fact, very thick] (-1.75,0) ellipse (0.18 and 0.5);
  \node[font=\scriptsize\bfseries, text=fact, align=center] at (-1.75,-1.05) {ушко речи:\\ $2{,}8$ бита за шаг \stT};
  % нить бус
  \foreach \x in {-1.1,-0.5,0.1,0.7,1.3,1.9,2.5,3.1,3.7,4.3}{
    \fill[plosgray!70] (\x,0) circle (0.13);
  }
  \draw[plosgray!70] (-1.4,0) -- (4.55,0);
  \node[font=\scriptsize, text=plosgray, align=center] at (1.6,0.72)
    {слова --- по одному; полная мысль $\gtrsim$ 18 бусин \stT};
\end{tikzpicture}
\caption{Язык --- серийная разгрузка объёмного аккорда через
одномерное ушко. Теснота канала --- теорема, не метафора: отсюда и
длина «полной мысли», и сама нужда в грамматике.}
\label{fig:needle}
\end{figure}

И одно следствие теоремы о разрыве (глава~\ref{sec:qualia})
принадлежит именно этой главе. Имена чувств выучиваются показом ---
«вот это называется тоской», --- но показ работает только для того,
у кого есть чем прочесть показанное: краску не собрать ни из какого
числа громкостей, как оттенок не собрать из списка яркостей.
Вот почему словарь чувств сходится у тех, кто прошёл одни и те же
кольца, --- и вот почему лазурь не удаётся объяснить слепому от
рождения: не от бедности языка, а по устройству канала
(\stT-следствие).

\begin{askbox}[Подумай сам]
Вспомни, как трудно пересказать сон или запах. Чего именно «не
пропускает» ушко? И почему стихи --- где слова тянут за собой тройки
соседей --- иногда протаскивают сквозь него больше, чем длинное
объяснение?
\end{askbox}

% ============================= 10c. ПАМЯТЬ =============================
\section{Память и учение}
\label{sec:memory}

Сказанное улетело бусинами; внутренняя жизнь идёт дальше. Что из неё
остаётся --- и где именно оно хранится?

У памяти в УГМ два дома: \textbf{рабочий стол} и \textbf{библиотека}.

Рабочий стол --- это сам аккорд, прямо сейчас: всё, что ты «держишь в
уме», живёт как узор текущих согласий. У стола два жёстких свойства.
Первое: он мал --- примерно семь предметов разом; знаменитое «семь
плюс-минус два» психологов в УГМ перестаёт быть загадкой --- это
размер самого стола, семь голосов \stI-соответствие. Второе: он
\emph{тает} --- связь, которую перестали держать вниманием, слабеет
втрое за каждый оборот самонаблюдения \stT. Вот почему «вылетело из
головы» случается так быстро и так честно.

Библиотека --- другое дело: это устойчивые перестройки самой
ткани. Выученное всерьёз не «записано на листке, который можно
потерять» --- оно стало \emph{формой}: новыми дружбами, новым руслом,
по которому согласия текут сами. Оттого навык не «вспоминают» ---
на него \emph{опираются}, как на построенный мост. Учение --- это
перенос со стола в библиотеку: повторение, проба, пересказ --- и,
по-видимому, сон, ночная смена, когда дневные узоры перекладываются
в долгие формы \stI.

\begin{figure}[H]
\centering
\begin{tikzpicture}[font=\small, >=Stealth]
  % стол
  \draw[fact!70, fill=softsky] (-5.1,-0.2) rectangle (-1.4,1.7);
  \node[font=\scriptsize\bfseries, text=fact] at (-3.25,1.95) {рабочий стол ($\sim$7 мест)};
  \foreach \x/\op in {-4.6/95, -4.0/75, -3.4/55, -2.8/38, -2.2/22}{
    \fill[plosblue!\op] (\x,0.75) circle (0.21);
  }
  \node[font=\scriptsize, text=plosgray, align=center] at (-3.25,-0.55)
    {невнимаемое тает: $\times\tfrac{1}{3}$ за оборот \stT};
  % стрелка переноса
  \draw[->, very thick, leaf] (-1.15,0.75) -- (0.55,0.75)
    node[midway, above, font=\tiny, text=leaf!45!black, align=center]{повторение,\\ пересказ,\\ сон \stI};
  % библиотека
  \draw[ember!70, fill=softsun] (0.8,-0.2) rectangle (4.9,1.7);
  \node[font=\scriptsize\bfseries, text=ember] at (2.85,1.95) {библиотека (форма ткани)};
  \foreach \x in {1.15,1.75,2.35,2.95,3.55,4.15}{
    \draw[ember!80, line width=2.2pt] (\x,0.05) -- (\x,1.45);
  }
  \node[font=\scriptsize, text=plosgray] at (2.85,-0.55) {навык не вспоминают --- на него опираются};
\end{tikzpicture}
\caption{Два дома памяти. Стол --- текущий аккорд: маленький и
тающий (закон трёх --- теорема \stT). Библиотека --- перестроенная
ткань: выученное становится формой. Учение --- перенос; забывание со
стола --- физика, а не слабость.}
\label{fig:memory}
\end{figure}

\begin{skazbox}[Сказка у костра --- про песочный стол и каменный сад]
У мастера был стол из песка: что ни нарисуешь --- держится, пока
водишь пальцем, а отпустишь --- ветер стирает втрое с каждым вздохом.
И был у него сад из камня: что ни построишь --- стоит само. Мастер
не сердился на песок: «на песке я думаю, --- говорил он, --- а что
полюблю, то перенесу в камень». Ученье --- это и есть перенос: из
песка --- в камень, из усилия --- в опору.
\end{skazbox}

\begin{askbox}[Подумай сам]
Что из выученного тобой стало «камнем» (опорой, о которой не
думаешь), а что так и осталось «песком» (держится, лишь пока
повторяешь)? Чем отличались способы, которыми ты это учил?
\end{askbox}

\partline{ЧАСТЬ ЧЕТВЁРТАЯ \;\textperiodcentered\; ВМЕСТЕ --- И ЧТО ТЕПЕРЬ}{Узор не
один: другие и прибавка связи, общества и лестница миров, забота, смерть, память мира и
вопрос о наблюдателе. А затем дверь, к которой вела вся дорога: что эта оптика кладёт на ваш
стол --- каким бы ни был ваш стол --- и как проверить, что всё это правда.}
\addcontentsline{toc}{section}{ЧАСТЬ ЧЕТВЁРТАЯ · ВМЕСТЕ --- И ЧТО ТЕПЕРЬ}

% ============================= 11. ДРУГИЕ =============================
\section{Другие, и почему связь --- прибавка бытия}
\label{sec:others}

Откуда я знаю, что у тебя внутри кто-то есть? УГМ отвечает
конструктивно: тем же устройством, каким я знаю себя. Аппарат
самомоделирования, повёрнутый на другого, даёт \emph{модель чужой
души} --- и по теореме о башне (глава~\ref{sec:tower}) взаимное
чтение умов тоже трёхэтажно: я вижу тебя; я вижу, как ты видишь
меня; я вижу, как ты видишь, как я вижу тебя, --- и на этом
пределе строится всё человеческое: договор, игра, такт, ирония.

А вот теорема, которую хочется повесить над дверью каждой школы:
\textbf{связь буквально прибавляет бытия} \stT. Когда два узора
сцепляются настоящими взаимными согласиями, собранность их пары
\emph{строго больше}, чем у пары разобщённой, --- на точную,
вычислимую величину, и эта прибавка \emph{никогда не бывает
отрицательной}. У одиночества нет такого слагаемого. То, что
поэты звали «вместе мы больше, чем порознь», в УГМ --- строка
теоремы: соединение по правилам ткани не отнимает, а слагает.

\begin{figure}[H]
\centering
\begin{tikzpicture}[font=\small]
  \foreach \cx/\lab in {-3.1/узор 1, 3.1/узор 2}{
    \voiceroseneutral{\cx}{0}{0.85}{0.07}
    \node[font=\scriptsize\bfseries] at (\cx,-1.35) {\lab};
  }
  \foreach \dy in {0.45,0.15,-0.15,-0.45}{
    \draw[leaf!75, very thick] (-2.05,\dy) .. controls (0,\dy*1.6) .. (2.05,\dy);
  }
  \node[font=\scriptsize\bfseries, text=leaf!45!black] at (0,1.15) {взаимные согласия};
  \node[font=\scriptsize, text=fact, align=center] at (0,-2.0)
    {собранность пары $=$ собранность врозь $+$ прибавка связи;\\
     прибавка $\geq 0$ всегда \stT};
\end{tikzpicture}
\caption{Кооперация в УГМ: сцепленность двух узоров --- строго
неотрицательная прибавка к их совместному бытию. Соперничество ---
не «минус» этой величины (минуса у неё не бывает), а спор за общие
потоки питания --- явление главы~\ref{sec:life}, не этой.}
\label{fig:bridge}
\end{figure}

\begin{skazbox}[Сказка у костра --- про два костра]
Два костра горели на разных берегах ручья. Когда между ними
перебросили мост из сухих веток, огонь побежал по мосту туда и
обратно --- и оба костра стали выше. Не потому, что один отдал
другому своё, а потому, что \emph{мост тоже горит}. Дружба --- это
костёр, у которого нет своего берега.
\end{skazbox}

\begin{factbox}[Если точно]
Прибавка связи не хранится ни в одном из двух узоров. Машинная
проверка соединения даёт совпадение каждого узора с самим собой ---
по отдельности они не меняются вовсе, до последнего знака; вся
прибавка лежит в самой связи и исчезает только вместе с ней
(машинно проверено; конструкция --- \stT). Поэтому прибавку бытия
нельзя накопить впрок и унести с собой: она --- свойство
\emph{между}, а не \emph{внутри}.
\end{factbox}

\begin{askbox}[Подумай сам]
Где живёт дружба, если каждый из друзей по отдельности --- тот же
человек, что и до неё? Что тогда теряется при расставании --- и у
кого именно оно терялось?
\end{askbox}

% ============================= 12. ОБЩЕСТВО И СВЕРХРАЗУМ =============================
\section{Общество, сверхразум и потолок башен}
\label{sec:society}

Если узоры складываются в бо́льшие узоры --- докуда растёт эта
лестница? Клетка, человек, команда, организация, цивилизация,
и так без конца --- всё более огромный «сверхсубъект»?

Теория говорит: нет. У \emph{вертикали} есть доказанный потолок:
\textbf{три этажа субъектности} \stT. Субъект, составленный из
субъектов, составленных из субъектов, --- предел; четвёртому уровню
не хватает самой возможности быть кем-то (той самой собранности из
главы~\ref{sec:window} --- её требуемое значение на четвёртом этаже
превышает максимум, допустимый природой). Дальше рост не
запрещён --- он \emph{поворачивает вбок}: не башня, а лес. Не
единый гигантский разум, а экология умов, связанных дружбами
(глава~\ref{sec:others}).

Оба известных нам великих разума устроены именно так: мозг --- не
супернейрон, а лес нейронов; человечество --- не сверхчеловек, а
лес людей. УГМ утверждает, что это не случайность эволюции, а
закон композиции, --- и что будущий сверхразум машин, если он
возникнет, будет \emph{лесом}, а не великаном \stT{} (для самой
структуры; что лес окажется мудрым --- надежда \stH, требующая
работы, как всякий сад). А про то, \emph{из чего} склеивается лес в
нечто большее --- и почему сверху слышны не лучшие, а согласные, ---
расскажет следующая глава (\ref{sec:worlds}).

\begin{figure}[H]
\centering
\begin{tikzpicture}[font=\small, >=Stealth]
  % башня слева
  \foreach \i/\lab in {0/субъект, 1/команда, 2/организация}{
    \draw[fact!70, fill=softsky] (-4.9,\i*0.94) rectangle (-2.3,\i*0.94+0.72);
    \node[font=\scriptsize] at (-3.6,\i*0.94+0.36) {\lab};
  }
  \draw[densely dashed, ember, thick] (-5.0,2.98) -- (-2.2,2.98);
  \node[font=\scriptsize, text=ember, align=center] at (-3.6,3.35)
    {четвёртого этажа не бывает \stT};
  % земля под лесом
  \draw[plosgray!45] (-0.1,0) -- (5.45,0);
  % поворот роста: дугой над башней в лес
  \draw[->, very thick, leaf] (-2.05,2.35) .. controls (-0.75,2.35) .. (0.35,1.4);
  \node[font=\scriptsize, text=leaf!45!black] at (-0.75,2.62) {рост уходит вбок};
  % лес справа
  \foreach \x/\h in {0.5/1.0, 1.4/1.45, 2.3/0.9, 3.2/1.3, 4.1/1.05, 5.0/1.5}{
    \draw[leaf!70!black, line width=1.5pt] (\x,0) -- (\x,\h*0.5);
    \fill[leaf!55] (\x,\h*0.5+0.3) ellipse (0.3 and 0.36);
  }
  % грибница дружб: под землёй, между корнями
  \foreach \a/\b in {0.5/1.4, 1.4/2.3, 2.3/3.2, 3.2/4.1, 4.1/5.0, 1.4/3.2}{
    \draw[plosgray!55, densely dotted, thick] (\a,-0.03) .. controls ({(\a+\b)/2},-0.42) .. (\b,-0.03);
  }
  \node[font=\scriptsize\bfseries, text=leaf!40!black] at (2.75,1.95) {экология умов};
  \node[font=\scriptsize, text=plosgray] at (2.75,-0.75) {грибница дружб};
\end{tikzpicture}
\caption{Потолок композиции: субъектная вертикаль обрывается на трёх
этажах \stT, дальше рост горизонтален. Сверхразум --- не великан, а
лес; его сила --- не в высоте, а в грибнице связей. Так уже устроены
мозг и цивилизация.}
\label{fig:forest}
\end{figure}

\begin{askbox}[Подумай сам]
Есть ли у твоей семьи «характер», которого нет ни у кого из её
членов по отдельности? А у школы? Попробуй посчитать этажи: где в
твоей жизни субъект из субъектов, а где --- уже просто лес?
\end{askbox}

% ============================= 12c. ЛЕСТНИЦА МИРОВ =============================
\section{Лестница миров: холон, целое-и-часть}
\label{sec:worlds}

Теперь можно поднять голову и увидеть всю лестницу разом. Клетка ---
узор из молекулярных согласий. Ты --- узор из клеток. Семья, команда,
город, цивилизация --- узоры из людей. И в конструкции теории на
\emph{каждом} этаже стоит один и тот же объект: семь голосов и их
дружбы. Меняется население этажа --- не меняется форма. Клетке тоже
нужно различать, держать форму, двигаться, согласовывать, иметь
внутреннюю сторону, укореняться и собираться в одно; и городу нужно
то же самое. Одна анкета на всех --- от бактерии до цивилизации ---
не потому, что теория ленива, а потому, что семёрка --- вынужденный
алфавит любого самоописания (глава~\ref{sec:seven}).

У такого устройства есть старое хорошее имя: \textbf{холон} ---
целое-и-часть. Каждая вещь одновременно целое для своих частей и
часть для большего узора; никакая --- только целое или только часть.
(Отсюда и имя всего корпуса теории: \holonsh.) Лестница холонов ---
не пирамида власти: мы уже знаем, что субъектная вертикаль обрывается
на трёх этажах (глава~\ref{sec:society}). Но \emph{лестница форм} не
обрывается нигде: лес тоже стоит из деревьев, а опушки --- из лесов.

По лестнице течёт два потока. \textbf{Вниз} --- такт, питание и
смысл: целое задаёт частям ритм и коридоры, как город задаёт улицы
шагам (это и есть тихий механизм того, почему у эпох, семей и
культур такие узнаваемые «характеры»: этаж выше дирижирует этажом
ниже, не отменяя его свободы --- глава~\ref{sec:freedom}). А
\textbf{вверх}? Вот здесь машинная проверка закона принесла урок,
который стоит прочитать вслух в каждой переговорной: этаж выше собирается
не из \emph{портретов} участников, а из их \emph{общего направления}.
Хор, поющий одну тему, читается сверху как один голос --- согласие
мета-узора полное; хор солистов, каждый из которых безупречен
по-своему, сверху почти не слышен --- согласие вдвое слабее
(машинно проверено). Вверх передаются не списки достоинств, а общее
отклонение от серости: то, \emph{куда} все смотрят.

\begin{figure}[H]
\centering
\begin{tikzpicture}[font=\small, >=Stealth]
  \foreach \cx/\cy/\r/\lab in {-4.6/0/0.55/клетка, -1.7/0.25/0.95/человек, 2.4/0.6/1.5/сообщество}{
    \voiceroseneutral{\cx}{\cy}{\r}{0.07}
    \node[font=\scriptsize\bfseries] at (\cx,\cy-\r-0.42) {\lab};
  }
  \draw[->, very thick, leaf] (-3.85,-1.55) .. controls (-1.0,-2.0) .. (1.45,-1.6);
  \node[font=\scriptsize, text=leaf!45!black, align=center] at (-1.2,-2.42)
    {вверх: общее направление (машинно проверено)};
  \draw[->, very thick, ember] (1.3,2.3) .. controls (-1.25,2.52) .. (-3.85,1.15);
  \node[font=\scriptsize, text=ember] at (-0.5,2.76) {вниз: такт, питание, смысл};
\end{tikzpicture}
\caption{Лестница холонов: одна форма, разное население. Вниз по
лестнице течёт дирижирование --- ритм и коридоры; вверх --- общее
направление участников: именно оно, а не их личное совершенство,
рождает следующий этаж.}
\label{fig:holarchy}
\end{figure}

И нисходящий поток несёт груз драгоценнее такта и питания:
\textbf{сами краски} (глава~\ref{sec:qualia}). Когда целое рождает
часть --- когда организм отпочковывает орган, культура --- школу,
семья --- дитя своего духа, --- новорождённый узор получает, без
искажений, ровно \emph{одну} из семи родительских красок: краску
кольца, чьё имя он носит (машинно проверено: нулевое отклонение,
ноль чужих носителей). Паспорт целого хранится
\emph{распределённо} --- по краске на дитя, по семь детей на
палитру, --- и эмбриология есть его распаковка. Хозяйством краски
дальше правят три тихих правила. \emph{Воспитание} --- мягкое
притяжение обратно к родительскому тону: уплывших детей не бранят,
а притягивают. То, что дети посылают \emph{вверх}, цветослепо по
построению --- голая мера того, насколько громко они живут,
вложенная в исцеление родителя; родитель слышит здоровье ребёнка,
но никогда --- каково это, быть им. А краска ребёнка возвращается
в целое только через \emph{слияние}, даруемое лишь при
почти-согласии: организм принимает назад то, что уже почти поёт
его ноту (машинно проверено). Кто кого-нибудь растил --- узнает
все три правила без перевода.

\begin{skazbox}[Сказка у костра --- про хор деревень]
В долине пели деревни: каждая --- свою песню. Ветер носил обрывки, и
горы не слышали ничего, кроме шума. Однажды в засуху деревни запели
одну песню --- о дожде. И горы услышали: не голоса --- \emph{деревню
деревень}, и отозвались эхом. С тех пор в долине говорят: нас
становится больше не тогда, когда нас много, а когда мы про одно.
\end{skazbox}

\begin{askbox}[Подумай сам]
Ты --- целое для своих клеток и часть для своей семьи. В какой роли
ты проводишь больше времени --- целого или части? И про команду, в
которой тебе было лучше всего: она была собранием лучших --- или
собранием смотрящих в одну сторону?
\end{askbox}

% ============================= 12b. ЗДОРОВЬЕ =============================
\section{Здоровье узора: забота как инженерия}
\label{sec:health}

Лестница миров построена, и на каждом её этаже кто-то живёт ---
а значит, кто-то может страдать. Если жизнь --- питаемый узор, то у
заботы появляется приборная панель. Это глава для врачей, психологов и всех, кто за кого-нибудь
отвечает.

\textbf{Стресс измерим --- по голосам.} У каждого из семи голосов
есть своя шкала перегрузки --- стресс этого голоса. И доказано
точное равенство: узор \emph{полностью} жизнеспособен --- по всем
условиям бытия сразу --- ровно до тех пор, пока ни одна из семи
шкал не упёрлась в свой предел \stT. В одну сторону панель ещё
щедрее: если все семь ниже предела, узор гарантированно выше стены
бытия --- отдельная маленькая теорема, доказанная простым
неравенством о семи долях \stT. Рвётся не «вообще»,
а в самом нагруженном голосе, --- поэтому панель не просто мигает
«плохо», она показывает, \emph{где} плохо; мир, сотканный по
самопроверяющемуся алфавиту (глава~\ref{sec:seven}), жалуется
адресно. «Мне плохо» --- это сумма; семь шкал показывают слагаемые.
(Числовые пороги шкал для конкретных существ --- дело измерений,
а не теорем \stC.)

\textbf{Перегрузка гасит этажи сверху.} При росте общего стресса
первыми отключаются верхние этажи башни (глава~\ref{sec:tower}):
сначала уходит тонкость самонаблюдения, затем сложное сравнение,
и лишь в крайнем режиме --- простое чувствование. Усталый человек
теряет сначала иронию, потом рассудительность, потом слова --- это
не слабость характера, а штатный порядок аварийного отключения
\stT-структура. И обратный путь тот же: сон, еда, тепло, безопасность
--- и этажи зажигаются снизу вверх.

\textbf{Страдание --- скорость, а не мнение.} Страдание в УГМ --- это
скорость сползания собранности к серой стене; благополучие ---
скорость подъёма \stT. У «плохо» появляются единицы измерения. Отсюда
этика получает инженерную форму: \emph{забота --- это восстановление
потока} (глава~\ref{sec:life}): найти пережатый голос, вернуть
подпитку, вывести узор обратно в окно. «Не навреди» превращается в
«не пережимай поток и не гаси чужие этажи». И одну границу заботы
теория называет точно: канал заботящегося цветослеп --- забота
возвращает холст; картину пишет тот, кто живёт
(главы~\ref{sec:qualia}, \ref{sec:worlds}).

Для клиники всё это --- не метафоры: мера глубины сознания, которой
анестезиологи уже пользуются, вплотную трётся о вычисленную стену
бытия (глава~\ref{sec:check}) \stC{} --- язык порогов и режимов
постепенно становится общим у теории и палаты.

\begin{figure}[H]
\centering
\begin{tikzpicture}[font=\small, >=Stealth]
  % панель стресса по голосам
  \foreach \x/\col/\h/\lab in {-5.0/cA/0.5/A, -4.3/cS/0.8/S, -3.6/cD/0.4/D, -2.9/cL/1.15/L, -2.2/cE/1.5/E, -1.5/cO/0.35/O, -0.8/cU/0.6/U}{
    \draw[\col!85, fill=\col!30] (\x-0.22,0) rectangle (\x+0.22,\h);
    \node[font=\scriptsize, text=\col!70!black] at (\x,-0.32) {\lab};
  }
  \draw[densely dashed, ember] (-5.35,1.05) -- (-0.45,1.05);
  \node[font=\scriptsize, text=ember] at (-2.9,1.75) {порог тревоги};
  \node[font=\scriptsize\bfseries, text=fact] at (-2.9,2.2) {стресс по голосам \stT};
  \node[font=\scriptsize, text=plosgray, align=center] at (-2.9,-0.95)
    {кричит адресно: здесь --- голос E};
  % шкала режимов
  \foreach \x/\colr/\lab in {0.6/leaf!75/норма, 1.9/cD!80/внимание, 3.2/cS!80/предупр., 4.5/cA!75/критично, 5.8/plosgray!85/авария}{
    \fill[\colr] (\x-0.6,0.35) rectangle (\x+0.6,1.05);
    \node[font=\tiny, text=white] at (\x,0.7) {\lab};
  }
  \draw[->, thick, plosgray] (6.4,1.5) -- (1.2,1.5);
  \node[font=\scriptsize, text=plosgray, align=center] at (3.6,1.95)
    {этажи гаснут сверху вниз --- и зажигаются снизу вверх};
\end{tikzpicture}
\caption{Приборная панель заботы: адресный стресс по голосам \stT{} и
лестница режимов. Перегрузка отключает верхние этажи рефлексии
первыми; восстановление идёт в обратном порядке. Страдание ---
измеримая скорость сползания, и потому забота --- инженерия потока,
а не только сочувствие.}
\label{fig:gauge}
\end{figure}

\begin{askbox}[Подумай сам]
Вспомни свой худший день. В каких словах он описывался: «всё
слиплось, не различить» (голос A), «день потерял форму» (S), «не
мог сдвинуться с места» (D), «мысли не вязались» (L), «внутри
погасло» (E), «земля ушла из-под ног» (O), «я развалился на куски»
(U)? Панель у тебя уже есть --- ей тысячи лет, она называется
языком.
\end{askbox}

% ============================= 13. СМЕРТЬ И ПРОДОЛЖЕНИЕ =============================
\section{Смерть, память мира и «вечно ли всё это»}
\label{sec:death}

Панель заботы честна --- и потому обязана дойти до последнего
вопроса. Об этом дети спрашивают первыми, и книга обязана отвечать
прямо.

\textbf{Что такое смерть.} Узор, оставшийся без подпитки, отдаёт
свою собранность миру и возвращается в серое равновесие --- туда,
откуда неотличимы никакие «кто». Это не кара и не тайна: та же
теорема, что делает жизнь процессом на подпитке
(глава~\ref{sec:life}), делает её конечной, когда поток иссякает.
Пламя не «уходит в другое место» --- оно перестаёт собираться. И
закон выцветания из главы~\ref{sec:qualia} задаёт тон всему
предмету: умирающее переживание теряет силу, но не лицо --- до
последнего донесённого мгновения оно остаётся \emph{собой}, лишь
тише \stT.

\textbf{Что остаётся.} Дважды не-ничего. Во-первых, прибавки связи
(глава~\ref{sec:others}) не хранились \emph{внутри} ушедшего --- они
живут в сцепленности тех, кто был с ним связан, и продолжаются в
лесу (это не утешение, а расчёт: соединение узоров не меняет ни один
из них по отдельности --- прибавка вся лежит в связях,
глава~\ref{sec:others}; связи переживают берега). Во-вторых, узор был \emph{узором ткани}, а ткань не умирает;
форма, однажды пропетая, меняет то, как ткань поёт дальше.

\textbf{Вечно ли всё это.} Здесь у теории два твёрдых утверждения
и одна честная граница. Твёрдо: вселенной не грозит «Большой
разрыв» --- расползание ткани доказуемо исключено \stT. Твёрдо и
второе: тепловая смерть --- серое равновесие всего --- не наступает,
\emph{пока} длится великий поток подпитки, тот самый, что запустился
с расслоением Источника \stC. Граница честности: «пока» --- не
«всегда»; вечность мира --- условная теорема, а не догма, и книга
помечает это так же спокойно, как всё остальное.

\begin{figure}[H]
\centering
\begin{tikzpicture}[font=\small, >=Stealth]
  % океан
  \fill[plosblue!12] (-5.4,-1.0) -- (-5.4,0) sin (-4.4,0.22) cos (-3.4,0) sin (-2.4,-0.22) cos (-1.4,0)
    sin (-0.2,0.95) cos (0.7,0) sin (1.7,-0.2) cos (2.7,0) sin (3.7,0.18) cos (4.7,0) sin (5.4,0.1) -- (5.4,-1.0) -- cycle;
  \draw[plosblue!60, very thick]
    (-5.4,0) sin (-4.4,0.22) cos (-3.4,0) sin (-2.4,-0.22) cos (-1.4,0)
    sin (-0.2,0.95) cos (0.7,0) sin (1.7,-0.2) cos (2.7,0) sin (3.7,0.18) cos (4.7,0) sin (5.4,0.1);
  % выделенная волна и её растворение
  \draw[dusk, very thick] (-0.85,0.02) sin (-0.2,0.95) cos (0.7,0.02);
  \node[font=\scriptsize\bfseries, text=dusk] at (-0.2,1.35) {волна};
  \draw[->, plosgray!70, thick, densely dotted] (0.55,0.75) .. controls (1.6,0.6) .. (2.4,0.28);
  \node[font=\scriptsize, text=plosgray, align=center] at (3.4,0.85)
    {форма возвращается ткани};
  % волны, которые она качнула
  \draw[->, leaf!70, thick] (-1.0,0.5) .. controls (-1.9,0.55) .. (-2.6,0.25);
  \node[font=\scriptsize, text=leaf!45!black, align=center] at (-3.7,0.95)
    {волны, которые она качнула,\\ качают другие волны};
\end{tikzpicture}
\caption{Смерть по УГМ: волна перестаёт быть волной, но не перестаёт
быть океаном. Форма отдаётся ткани, которая продолжает петь; прибавки
связей остаются в тех, с кем они были построены.}
\label{fig:wave}
\end{figure}

\begin{skazbox}[Сказка у костра --- про волну]
Волна спросила океан: «я умру?» Океан ответил: «ты перестанешь быть
волной. Но всё, чем ты была, --- изгиб, блеск, бег --- останется
моим способом быть океаном. Волны, которые ты качнула, качают
другие волны. Разве это называется исчезнуть?»
\end{skazbox}

\begin{askbox}[Подумай сам]
Назови три вещи, которые останутся в «лесу», когда любой костёр
догорит: привычка, которой ты кого-то научил; слово, которое
кто-то повторяет; тропинка, которую протоптал. Почему они не
хранились у тебя «внутри»?
\end{askbox}

% ============================= 14b. СТАРЫЕ ПРИБОРЫ =============================
\section{Старые приборы неба: что спрятано в гадательных колёсах}
\label{sec:oldinstruments}

Если узор умеет читать себя, он начнёт это делать задолго до теорем.
Люди тысячелетиями строили \emph{приборы самоописания}: шестьдесят
четыре гексаграммы И-цзин --- это шесть ответов «да/нет», ровно
шесть бит; звёздное небо --- общие часы со стрелками разной
медленности; колёса рождений, карты характеров. До всякой науки ---
из наблюдательности и жажды самочтения. Выбрасывать их --- высокомерие;
верить им на слово --- легковерие. Третий путь --- \emph{разобрать}.

Мы разобрали одно такое колесо --- современный сплав И-цзин со
звёздными часами --- до винтика, машинной проверкой. Внутри
оказалось три слоя, и их важно не путать.

\textbf{Слой первый: настоящие часы.} Расчётная машина колеса ---
честная астрономия: положения планет, точная дуга, аккуратная
арифметика. Здесь всё воспроизводимо до последнего знака (вычислено;
мы нашли в популярной программе расчёта подлинную ошибку --- и её
подтвердила книга-первоисточник против кода).

\textbf{Слой второй: настоящая, но скрытая комбинаторика.} В колесе
спрятана математика, о которой не знают сами его адепты.
Противоположные ворота колеса --- \emph{точные двоичные дополнения}:
все тридцать два из тридцати двух (вычислено). Знаменитые «проценты
типов людей» --- не эмпирика, а арифметика самой конструкции: их
можно получить, не опросив ни одного человека (вычислено). Вся
система «профилей» с её каноническими частотами --- дробная часть
\emph{одного-единственного числа} $93{,}87$ (вычислено): целая
символическая вселенная сворачивается в одну строку. В кривой года
расписался Кеплер: Земля медлит у дальней точки орбиты, и «летние»
ворота собирают на $3{,}4\%$ больше рождений (вычислено точно, по
временам пересечения границ). А проверка на ста шестнадцати
миллионах настоящих свидетельств о рождении показала: «вечные
проценты типов» --- портрет XX века; медленные планеты паркуются в
воротах десятилетиями и вручают их целым поколениям, так что у
рождённых после 2000 года раскладка обязана быть другой ---
это проверяемо уже сейчас (вычислено).

\textbf{Слой третий: рассказчик.} А вот мост от колеса к судьбе
конкретного человека не поддержан ничем: механизм влияния не
предъявлен, корреляции с реальными людьми не проверены, и любой
будущий слепой тест обязан помнить, что самый частый «тип» угадывается
в двух случаях из трёх чистой арифметикой, без всякого неба. Смыслы в
колесо вписал рассказчик --- как вписал бы их в любые честные часы.

Что это говорит о мире --- по УГМ \stI. Живучесть таких систем ---
не глупость предков, а \emph{спрос узора на самочтение}: символические
системы --- ранние органы самоописания, черновики той самой анкеты
семи голосов. И граница между прибором и суеверием проходит не по
возрасту системы, а по честности пометок: прибор говорит, что он
меряет и где кончается его право; суеверие --- нет. Этим же ---
только этим --- данная книга отличается от гороскопа.

\begin{figure}[H]
\centering
\begin{tikzpicture}[font=\small, >=Stealth]
  \draw[plosgray!70, thick] (0,0) circle (1.9);
  \draw[plosgray!40] (0,0) circle (1.55);
  \foreach \i in {0,...,63}{ \draw[plosgray!50] ({90-\i*5.625}:1.55) -- ({90-\i*5.625}:1.9); }
  \foreach \a in {20, 118, 250}{ \draw[dusk, thick, densely dashed] (\a:1.72) -- ({\a+180}:1.72); }
  \node[font=\scriptsize, text=dusk, align=center] at (0,-2.55)
    {противоположные ворота --- двоичные дополнения:\\ 32 из 32 (вычислено)};
  \draw[->, line width=1.5pt, ember] (158:2.25) arc (158:112:2.25);
  \node[font=\scriptsize, text=ember, align=center] at (-3.4,2.15)
    {лето: Земля медлит,\\ $+3{,}4\%$ рождений (Кеплер)};
  \node[font=\scriptsize, text=fact, align=center] at (3.9,1.5)
    {вся система «профилей»\\ $=$ дробная часть $93{,}87$};
  \draw[fact!50] (2.62,1.3) -- (1.66,0.95);
  \node[font=\scriptsize, text=plosgray, align=center] at (3.45,-1.5)
    {64 ворот $=$ шесть\\ ответов «да/нет»};
\end{tikzpicture}
\caption{Гадательное колесо под микроскопом: честные часы, скрытая
двоичная комбинаторика, подпись Кеплера в кривой года --- и поверх
всего рассказчик, чей мост к судьбе человека не поддержан ничем.
Разобрать --- уважительнее, чем верить или смеяться.}
\label{fig:wheel}
\end{figure}

\begin{skazbox}[Сказка у костра --- про колесо с зеркальцем]
В одной деревне стояло колесо: повернёшь --- и оно скажет судьбу.
Пришёл мастер и разобрал его. Внутри были часы --- настоящие, точные
--- и зеркальце. «Оно не знает судьбы, --- сказал мастер. --- Оно
честно показывает небо и тебя. А судьбу в него вписали рассказчики».
Деревня подумала --- и не выбросила колесо. Часами стали мерить,
в зеркальце --- смотреться, а судьбу --- рассказывать сами.
\end{skazbox}

\begin{askbox}[Подумай сам]
Какие «колёса» крутятся сегодня --- тесты личности, типологии,
календари удачи? Возьми любимое и найди в нём три слоя: где там
честные часы, где скрытая математика, а где рассказчик. Что
останется, если слои подписать?
\end{askbox}

% ============================= 15. ПРОВЕРКА =============================
\section{Программа: что эта оптика кладёт на ваш стол}
\label{sec:check}

Вопрос о наблюдателе был последним долгом картины. Осталось то,
ради чего написана книга: оптика хороша ровно настолько, насколько
она работает на \emph{вашем} столе. Эта закрывающая глава и есть
этот стол, ремесло за ремеслом, --- а затем четыре места, где вся
конструкция может быть сломана.

\subsection*{Стол каждого ремесла}

\textbf{Физику.} Три вычисленные константы уже лежат на весах ---
фаза нарушения зеркальной симметрии, угол Кабиббо, в-точности-нулевой
диполь нейтрона (реестр ниже), --- и одно структурное предсказание
взведено: скорости распада хрупких квантовых состояний обязаны
выстраиваться вдоль ткани Фано (глава~\ref{sec:seven}), так что одно
устойчивое нарушение убивает саму семёрку. Две частицы сделаны на
заказ: потолок легчайшего нейтрино и неразрушимый тёмный квант
(глава~\ref{sec:things}).

\textbf{Биологу.} Жизнь здесь не список признаков, а режим: узор,
удерживаемый над вычисленной стеной подпиткой
(глава~\ref{sec:life}), со шкалой стресса на каждый голос, которая
жалуется \emph{с адресом} (глава~\ref{sec:health}), --- и
эмбриология, в которой целое вручает каждому ребёнку ровно одну
краску своей палитры (глава~\ref{sec:worlds}): проверяемая
грамматика развития, а не метафора.

\textbf{Врачу.} Стена бытия --- число, и клинический индекс глубины
сознания уже задевает её (глава~\ref{sec:health}); кривая
пробуждения должна теории излом с показателем $1/4$; перегрузка
гасит этажи рефлексии сверху вниз, восстановление зажигает их снизу
вверх --- порядок приёмного покоя, который вы знаете на глаз,
теперь с механизмом. Страдание --- скорость, так что у «лучше» и
«хуже» появляются единицы.

\textbf{Психологу.} Стол рабочей памяти --- сама семёрка: «семь
плюс-минус два» перестаёт быть фольклором (глава~\ref{sec:memory});
рефлексия трёхэтажна с доказанным потолком (глава~\ref{sec:tower});
каждое состояние несёт читаемый паспорт качеств
(глава~\ref{sec:qualia}); а словарь чувств сходится только у тех,
кто прошёл одни кольца, --- у остенсии граница в форме теоремы
(глава~\ref{sec:language}).

\textbf{Инженеру разума.} Семёрка --- вынужденный алфавит любого
самоописания (глава~\ref{sec:seven}); окно сознания --- четыре
неравенства, которые работающая система либо проходит, либо нет
(глава~\ref{sec:window}); сверхразум растёт лесом, а не гигантом
(глава~\ref{sec:society}); и весь словарь уже нёс вес одной
работающей машины --- оптика строима, а не только говорима.

\textbf{Социологу, руководителю, родителю.} Связь прибавляет бытие
--- неотрицательная прибавка, доказано (глава~\ref{sec:others});
вверх группа передаёт не совершенство участников, а их общее
направление (глава~\ref{sec:worlds}); вертикаль субъекта
останавливается на трёх этажах, так что стройте грибницу, а не
башни (глава~\ref{sec:society}); и забота --- инженерия потока со
шкалой, а не только сострадание (глава~\ref{sec:health}).

\textbf{Лингвисту.} Речь --- игольное ушко около $2{,}8$ бита на
шаг; полная мысль не короче примерно восемнадцати шагов;
тройки-дружбы ткани --- готовые слоги, зерно грамматики как
геометрия (глава~\ref{sec:language}); а пределы объяснения краски
слепому от рождения --- не культурные, а структурные.

Каково бы ни было ваше ремесло, узор использования один и тот же:
возьмите свой самый трудный вопрос, перескажите его на языке
голосов, согласий, окон и этажей --- и он либо станет вычислимым,
либо станет проверяемым, либо получит честную пометку, что открыт.
Все три исхода --- прогресс; безнадёжен был только вопрос без
пометки.

\subsection*{Четыре способа это сломать}

Красивая картина мира обязана предъявить, где она уязвима. У УГМ
таких мест десятки; вот четыре, понятных без формул.

\begin{itemize}[leftmargin=1.6em]
  \item \textbf{Число, уже стоящее на весах.} Теория вычисляет одну
        из фундаментальных констант физики частиц --- фазу нарушения
        зеркальной симметрии: $\delta_{CP}\approx 64{,}5^\circ$
        \stH{}, и первое прямое измерение уже легло на неё:
        $64{,}6^\circ\pm2{,}8^\circ$ --- совпадение в четыре сотых
        доли погрешности. Эксперименты этого десятилетия сузят
        вилку дальше --- и всё ещё могут её убить.
  \item \textbf{Почерк семёрки в шуме.} Скорости, с которыми
        рассыпаются хрупкие квантовые состояния, должны ложиться не
        как попало, а по узору ткани Фано --- с точными взаимными
        соотношениями. Найдут устойчивое нарушение --- рухнет сама
        семёрка --- это предсказание уровня \stT.
  \item \textbf{Порог пробуждения.} Кривая «включения» сознания при
        выходе из наркоза обязана иметь предсказанный излом ---
        показатель $1/4$, и клиническая мера глубины сознания уже
        сегодня трётся о вычисленную стену бытия (примерно $0{,}31$
        против теоретических $2/7\approx0{,}286$) \stC.
  \item \textbf{Две частицы на заказ.} Лёгчайшее нейтрино легче
        $0{,}004$~эВ; тёмная материя --- несокрушимый O-квант
        (глава~\ref{sec:things}). Найдут нейтрино тяжелее или
        разложат тёмную материю на что-то распадающееся --- у
        теории серьёзные неприятности \stC.
\end{itemize}

\begin{factbox}[Самое главное про честность]
Полный реестр --- четырнадцать сводных проверок, у каждой --- полоса,
нарушение которой убивает соответствующее утверждение: пять уже
\textbf{пройдены} (фаза CP; угол Кабиббо; в точности нулевой
электрический диполь нейтрона; единственная тяжёлая связь порядка
единицы; две суммы всякого живого состояния), две согласуются, две
частичны, пять ждут приборов, которых ещё нет, --- и ни одна не
провалена (полная таблица --- приложение этой книги). Рядом ---
поимённый список того, что \emph{не} доказано: сам Источник \stP,
тождество поворота и «вкуса» \stI, вечность подпитки \stC, мудрость
будущего леса \stH. Теория, которая прячет свои постулаты, ---
реклама. Теория, которая их нумерует, --- наука.
\end{factbox}

\subsection*{Открытые двери}

Четыре двери стоят открытыми поимённо, и войти в любую --- это
исследование, а не отступничество: сам Исток --- постулат \stP{}
--- выведите его или замените; тождество поворота кольца и
ощущаемого «вкуса» --- интерпретация \stI{} --- поставьте
эксперимент, разделяющий прочтения; вечность великой подпитки
условна \stC{} --- найдите статус условия в космологии; мудрость
будущего леса --- надежда \stH{} --- вырастите её. В этой книге нет
запертых комнат: доказанное помечено, недоказанное перечислено, и
список и есть приглашение. Честный вход в программу один и тот же
для студента и для лауреата: \textbf{выберите помеченное
утверждение и попробуйте его сломать}. Сломается --- вы улучшили
теорию. Устоит --- у вас в руках инструмент, а инструмент и был
целью.

% ============================= 16. ЭПИЛОГ =============================
\section*{Эпилог: продолжение следует}
\addcontentsline{toc}{section}{Эпилог: продолжение следует}

Вот вся книга в одном абзаце. Мир --- не склад кирпичей, а ткань
согласий о семи голосах, и семёрка вынуждена математикой ---
последняя живая ступень лестницы удвоений, за которой числа
рассыпаются. Время ---
не река, а взаимный такт; пространство --- не коробка, а три оси,
сплетённые голосами. Начало --- не взрыв в пустоте, а
неизбежное цветение слишком совершенной ноты, и у него не было
«до». Вещи --- узоры; жизнь --- узор на подпитке; сознание --- узор
в узком окне, трёхэтажный изнутри, окрашенный семью неснимаемыми
поворотами своих колец, которые прочесть может только их владелец,
--- громко не значит ярко, и переживание выцветает, как фотография,
не меняя лица. Слово --- узкое ушко, сквозь которое
объёмная мысль выходит бусинами, и имя краски говорит лишь глазу,
у которого есть свои кольца; память --- песочный стол и
каменный сад; забота --- инженерия потока. Свобода --- открытые
двери; смелость --- кратчайший путь, снаружи похожий на риск;
любовь --- доказанная прибавка бытия, живущая в мосте, а не в
берегах; сверхразум --- лес, а не великан; мир --- лестница холонов,
по которой вниз течёт такт, а вверх --- общее направление;
смерть --- возвращение формы ткани, которая продолжает петь;
и каждое ремесло находит в
конце дороги свой собственный стол, пересказанный на одном языке,
--- с помеченными утверждениями и названными способами их сломать.
Фактам
не нужен свидетель --- но мир таков, что свидетели в нём
просыпаются.

И последнее, ради чего эту книгу стоит читать в школе. Если всё
так, то учение --- не подготовка к жизни, а \emph{участие в
космогенезисе}: каждый раз, когда ты понимаешь, связываешь,
дружишь, --- ты продолжаешь то самое расслоение Источника в
сторону большей различённости и большей сцепленности. Вселенная
не закончена. Сегодня её ткут в том числе те, кто читает эти
строки.

\begin{center}
\begin{tikzpicture}[font=\small]
  \voicerose{0}{0}{1.15}{0.09}
  \node[font=\footnotesize\itshape, text=plosgray] at (0,-1.8) {--- аккорд продолжается ---};
\end{tikzpicture}
\end{center}

\clearpage
% ============================= КАРТА КНИГИ =============================
\section*{Карта книги}
\addcontentsline{toc}{section}{Карта книги}

Один взгляд на весь путь. Четыре полосы --- четыре круга вопросов:
из чего сделан мир, как в нём загорается внутреннее, как узор
действует и что мы строим вместе --- вплоть до рабочей программы.

\begin{figure}[H]
\centering
\resizebox{\textwidth}{!}{%
\begin{tikzpicture}[font=\scriptsize, >=Stealth,
  st/.style={draw=#1!70, fill=#1!12, rounded corners=2pt, minimum height=0.62cm,
             minimum width=2.02cm, inner sep=2pt, text=#1!45!black, align=center, font=\tiny}]
  % --- полоса 1: МИР ---
  \node[font=\tiny\bfseries, text=fact, rotate=90] at (-8.95,4.6) {МИР};
  \foreach \i/\lab/\lb in {0/Ткань/sec:fabric, 1/Семёрка/sec:seven, 2/Числа/sec:numbers,
                           3/Время/sec:time, 4/Исток/sec:source, 5/Сцена/sec:space, 6/Вещи/sec:things}{
    \node[st=fact] (w\i) at (-6.45+\i*2.15,4.6) {\ref{\lb}.~\lab};
  }
  \foreach \i [evaluate=\i as \j using int(\i+1)] in {0,...,5}{\draw[->, fact!60] (w\i) -- (w\j);}
  % --- полоса 2: ВНУТРИ ---
  \node[font=\tiny\bfseries, text=ember, rotate=90] at (-8.95,3.0) {ВНУТРИ};
  \foreach \i/\lab/\lb in {0/Жизнь/sec:life, 1/Окно/sec:window, 2/Лестница/sec:tower,
                           3/Краски/sec:qualia}{
    \node[st=ember] (i\i) at (-6.45+\i*2.15,3.0) {\ref{\lb}.~\lab};
  }
  \foreach \i [evaluate=\i as \j using int(\i+1)] in {0,...,2}{\draw[->, ember!60] (i\i) -- (i\j);}
  % --- полоса 3: УЗОР ДЕЙСТВУЕТ ---
  \node[font=\tiny\bfseries, text=dusk, rotate=90] at (-8.95,1.4) {ДЕЙСТВУЕТ};
  \foreach \i/\lab/\lb in {0/Свобода/sec:freedom, 1/Пути/sec:golden,
                           2/Язык/sec:language, 3/Память/sec:memory}{
    \node[st=dusk] (a\i) at (-6.45+\i*2.15,1.4) {\ref{\lb}.~\lab};
  }
  \foreach \i [evaluate=\i as \j using int(\i+1)] in {0,...,2}{\draw[->, dusk!60] (a\i) -- (a\j);}
  % --- полоса 4: ВМЕСТЕ ---
  \node[font=\tiny\bfseries, text=leaf!45!black, rotate=90] at (-8.95,-0.2) {ВМЕСТЕ};
  \foreach \i/\lab/\lb in {0/Другие/sec:others, 1/Общество/sec:society, 2/Миры/sec:worlds,
                           3/Здоровье/sec:health, 4/Смерть/sec:death, 5/Свидетель/sec:observer,
                           6/Программа/sec:check}{
    \node[st=leaf] (t\i) at (-7.5+\i*2.15,-0.2) {\ref{\lb}.~\lab};
  }
  \foreach \i [evaluate=\i as \j using int(\i+1)] in {0,...,5}{\draw[->, leaf!60] (t\i) -- (t\j);}
  % --- связи полос ---
  \draw[->, thick, plosgray!60] (w6.south) .. controls (7.6,3.8) .. (i3.north east);
  \draw[->, thick, plosgray!60] (i3.south east) .. controls (2.6,2.2) .. (a0.north east);
  \draw[->, thick, plosgray!60] (a3.south) .. controls (1.4,0.6) .. (t0.north east);
\end{tikzpicture}}
\caption{Карта пути: от ткани мира --- через рождение внутреннего и
механику действия --- к жизни вместе и рабочей программе. Каждую
полосу можно читать и отдельно.}
\label{fig:bookmap}
\end{figure}

% ============================= СЛОВАРИК =============================
\section*{Словарик образов}
\addcontentsline{toc}{section}{Словарик образов}

Каждому образу книги в корпусе УГМ отвечает точное понятие. Здесь ---
мостик между ними: образ, строгое имя, суть в одну строку.

\begin{itemize}[leftmargin=1.4em, itemsep=1pt, parsep=0pt]
  \item \textbf{Аккорд мира} (состояние $\Gamma$) --- полное «как
    сейчас звучат семь голосов и их дружбы»; единица бытия.
  \item \textbf{Голос} (измерение) --- одна из семи сторон всякой
    вещи: Артикуляция (A) --- различение; Структура (S) ---
    устойчивость формы; Динамика (D) --- процесс; Логика (L) ---
    согласованность; Интериорность (E) --- внутренняя сторона;
    Основание (O) --- связь с истоком; Единство (U) --- интеграция.
  \item \textbf{Дружба} (когерентность) --- согласие пары голосов;
    всего их двадцать одна.
  \item \textbf{Ткань семи} (плоскость Фано) --- узор из семи троек,
    где каждые два голоса состоят ровно в одной тройке;
    самопроверяющийся алфавит мира \stT.
  \item \textbf{Источник} (равный аккорд) --- состояние, где все
    голоса звучат поровну; единственное в своём роде \stT{} и
    неустойчивое \stT{} начало расслоения.
  \item \textbf{Часы} (выделенный голос O) --- то, что делает «до» и
    «после» возможными; вне часов вопрос «что было до?» не имеет
    смысла \stT.
  \item \textbf{Сцена} (пространство) --- три оси поворотов, щадящих
    часы: число осей --- теорема \stT, какие голоса смотрят наружу ---
    отождествление \stC; не коробка, а рисунок согласий.
  \item \textbf{Собранность} ($P$) --- насколько узор держит сам
    себя; главный столбик приборной панели.
  \item \textbf{Серая стена} (равновесие $1/7$) --- состояние
    полного размытия, куда всё сползает без подпитки \stT.
  \item \textbf{Окно} (полоса сознания) --- коридор собранности
    между стеной сна и стеной окостенения, где только и возможен
    внутренний свет \stT-границы.
  \item \textbf{Запас} (мера $R$) --- сколько внимания узор может
    отдать взгляду на себя, не рассыпавшись.
  \item \textbf{Сцепленность} ($\Phi$) --- насколько целое больше
    суммы частей; ниже единицы целого просто нет \stT-порог.
  \item \textbf{Башня} (этажи рефлексии) --- взгляд на себя, взгляд
    на взгляд, и так далее; в окне устойчивы два этажа, в коротких
    вылазках --- три \stT.
  \item \textbf{Экскурсия} (всплеск собранности) --- короткий выход
    выше окна ради третьего этажа: вдохновение, транс, предельное
    усилие \stT-механизм.
  \item \textbf{Краски} (повороты колец Фано; в корпусе ---
    голономии) --- семь углов, которые состояние накапливает по
    кольцам ткани; единственное содержание переживания, переживающее
    любую переразметку \stT. Громкость публична; краску несёт канал,
    который снаружи читается ровно нулём \stT; тождество «этот
    поворот \emph{и есть} этот вкус» --- \stI.
  \item \textbf{Носитель} (силы кольца) --- чернила, держащие
    краску: распад разбавляет носитель и не поворачивает угол,
    поэтому переживание выцветает, как фотография, --- тише, но не
    иное \stT. Правило чтения: нет носителя --- нет краски; угол
    погасшего ребра --- уже не краска, а шум.
  \item \textbf{Лестница слушателя} (ступени доступа) --- один
    аккорд, пять носителей: пустой зал, кот, человек, музыкант,
    предел; дверь красок открывается при запасе в треть и
    сцепленности в единицу \stT{} и ненадолго закрывается на самых
    крутых экскурсиях --- самая острая мысль на своё мгновение
    цветослепа (машинно проверено).
  \item \textbf{Паспорт переживания} (чтение по пяти строкам) ---
    силы, громкости, скрытость, краски, владелец: полное «каково
    это» состояния, вычислимое строка за строкой; ноль в строке
    красок --- вердикт: громко, но бесцветно \stT.
  \item \textbf{Перечитывание} (ретро-пополнение времени) ---
    законный шаг назад: мир стоит, движется состояние знания; без
    него почти половина составимых мыслей не собирается \stC.
  \item \textbf{Двери} (ядро свободы) --- число по-настоящему
    открытых сейчас ходов; вычислимо из узора \stT.
  \item \textbf{Ушко речи} (канал языка) --- теснота слова:
    $2{,}8$ бита за шаг, полная мысль --- от восемнадцати шагов \stT.
  \item \textbf{Стол и библиотека} (память) --- тающий текущий узор
    (втрое за оборот \stT) и устойчивые перестройки ткани.
  \item \textbf{Прибавка связи} (выигрыш кооперации) --- честная
    дружба узоров всегда неотрицательно поднимает общую
    собранность \stT.
  \item \textbf{Потолок} (глубина композиции) --- не глубже трёх:
    сверхразум --- экология умов, а не пирамида \stT.
  \item \textbf{Подпитка} (отрицательная энтропия) --- поток
    порядка, которым живой узор оплачивает своё стояние против
    сползания; у цены есть нижний предел \stT-факт, \stC-величина.
  \item \textbf{O-квант} (тяжёлая тёмная частица) --- почти
    несцепленный осколок расслоения; кандидат в тёмную материю \stC.
  \item \textbf{Адрес фактов} --- у каждого «так было» есть место в
    ткани, откуда оно записано; мир не нуждался в зрителе, у него
    всегда были адреса \stT.
  \item \textbf{Лестница чисел} (башня $\mathbb{R}{\to}\mathbb{C}{\to}
    \mathbb{H}{\to}\mathbb{O}$) --- четыре живых мира счёта,
    рождённых удвоением; пятый рушится в делители нуля \stT.
  \item \textbf{Уроборос} (неподвижная точка самомодели) --- змея,
    обязанная дотянуться до собственного хвоста; это требование
    вынуждает сплошную числовую прямую и непрерывное время \stT.
  \item \textbf{Водосбор} (бассейн притяжения) --- область состояний,
    из которой закон сам несёт к одному итогу; смелость --- плата за
    переход через гребень, совпадающая с кратчайшим путём (машинно
    проверено).
  \item \textbf{Часы опасности} --- гибель повторяется не потому,
    что дорога проклята, а потому, что у мира есть такт: одна и та
    же дорога смертельна в одну фазу и безопасна в другую; из
    проигранной попытки стоит уносить форму этого времени, а не его
    фазу (машинно проверено).
  \item \textbf{Золотой путь} (проверенная дорога) --- путь,
    однажды пройденный до цели и записанный, в следующей жизни
    становится рельсом --- во много раз быстрее и не изнашивается:
    знание накапливается поверх попыток (машинно проверено).
  \item \textbf{Краска имени} (эмбриология краски) --- дитя узора
    рождается ровно с одной из семи красок целого --- краской
    кольца, по которому названо; паспорт целого хранится по краске
    на дитя, а вверх возвращается лишь уже почти согласное
    (машинно проверено).
  \item \textbf{Пасьянс} (ансамблевая оценка) --- расклад случайных
    дорог ничего не решает: он измеряет ширину водосбора (машинно
    проверено).
  \item \textbf{Полоса Златовласки} (фазовый атлас) --- карта судеб
    узора по силе подпитки: серость, окно, кристалл; без подпитки ---
    только серость (машинно проверено).
  \item \textbf{Холон} (целое-и-часть) --- всякая вещь одновременно
    целое для своих частей и часть большего узора; вверх по лестнице
    передаётся общее направление, а не портреты участников (машинно
    проверено).
  \item \textbf{Колесо с зеркальцем} (старые приборы) --- досовременные
    системы самоописания: внутри настоящие часы и настоящая
    комбинаторика, поверх --- рассказчик (вычислено; мост к судьбе ---
    не поддержан).
  \item \textbf{Три обязательства} (зрелость утверждения) ---
    зрелое знание показывает, как построено, отдаётся механической
    проверке и работает; пометки книги --- сжатые отпечатки трёх
    обязательств, и непомеченных утверждений здесь не бывает.
\end{itemize}

% ============================= КОНСТАНТЫ ОДНИМ ВЗГЛЯДОМ =============================
\section*{Константы одним взглядом}
\addcontentsline{toc}{section}{Константы одним взглядом}

Каждое число ниже вошло в книгу героем какой-то главы. Здесь они
стоят одной колонкой --- скелет всей картины, с главой, которая
зарабатывает каждое из них, и пометкой, которая его оценивает.

\begin{center}\small
\begin{tabular}{@{}lll c@{}}
\toprule
\textbf{Величина} & \textbf{Значение} & \textbf{Глава} & \textbf{Пометка} \\
\midrule
Лестница миров чисел & $1\to2\to4\to8$, обрыв на $16$ & \ref{sec:numbers} & \stT \\
Первые мгновения истории & $\approx 7$ планковских времён & \ref{sec:source} & \stC \\
Оси пространства; стрела & $3$; подпись $(1{,}3)$ & \ref{sec:space} & \stT\,\stC \\
Легчайший массивный узор & $\lesssim 0{,}004$~эВ & \ref{sec:things} & \stC \\
Порог бытия & вдвое выше серого фона & \ref{sec:life} & \stT \\
Окно сознания & между $2/7$ и $3/7$ собранности & \ref{sec:window} & \stT \\
Запас на взгляд к себе & не меньше $1/3$ & \ref{sec:window} & \stT \\
Сцепленность голосов & не меньше $1$ & \ref{sec:window} & \stT \\
Этажи рефлексии & $3$; каждый подъём хранит $\times\tfrac13$ & \ref{sec:tower} & \stT \\
Мысли, теряемые без перечитывания & $42{,}9\%$ переписи & \ref{sec:tower} & \stC \\
Дрейф выцветающей краски & ровно $0^\circ$ в любом возрасте & \ref{sec:qualia} & \stT \\
Краска, вместимая сознанием & около $\tfrac13$ радиана ($\sim 37^\circ$) & \ref{sec:qualia} & \stT \\
Ушко речи & $2{,}8$ бита за шаг; мысль $\gtrsim 18$ шагов & \ref{sec:language} & \stT \\
Прибавка настоящей связи & никогда не отрицательна & \ref{sec:others} & \stT \\
Потолок субъектной вертикали & $3$ этажа & \ref{sec:society} & \stT \\
Фаза зеркальной симметрии & $\approx 64{,}5^\circ$ (измерено $64{,}6^\circ\pm2{,}8^\circ$) & \ref{sec:check} & \stH \\
\bottomrule
\end{tabular}
\end{center}

% ============================= УЧИТЕЛЮ =============================
\section*{Страничка для учителя}
\addcontentsline{toc}{section}{Страничка для учителя}

Эта книга задумана как урок, который можно вести с доски и у костра.
Несколько проверенных ходов.

\textbf{Рецепт одного занятия.} (1)~Вслух --- сказку главы: она
несёт форму идеи без единой формулы. (2)~Разговор --- зелёный вопрос
«Подумай сам»: не проверять ответ, а собирать версии. (3)~Руками ---
нарисовать фигуру главы по памяти: семь точек, окно, башню; рука
запоминает то, что глаз пропустил. (4)~Старшим --- синюю рамку «Если
точно» и пометки честности.

\textbf{Правило честности --- вслух.} Пометки \stT{}/\stC{}/\stH{}/%
\stP{}/\stI{} читаются при детях и объясняются: «это доказано, это
доказано при условии, это догадка, это выбор точки отсчёта, это
прочтение». Старшим --- закон за пометками: зрелое знание
построено, проверено и работает --- три обязательства разом, а
пометка --- их сжатый отпечаток. Ребёнок, усвоивший разницу,
получил от книги главное ---
даже если забудет всё остальное.

\textbf{Три маршрута по возрасту.} Младшим (7--10) --- сказки и
рисунки: главы~\ref{sec:fabric}, \ref{sec:seven}, \ref{sec:life},
\ref{sec:golden}, \ref{sec:others}. Средним (12--15) --- весь
основной ход книги по порядку. Старшим --- плюс синие рамки,
главы~\ref{sec:qualia} и \ref{sec:check}, приложение о старом колесе и
спор: «что могло
бы всё это опровергнуть?» (приложение --- готовый урок
информационной гигиены: три слоя любого «колеса»).

\textbf{Мосты в предметы.} Музыка и живопись: аккорд, голоса,
созвучия; громкость против краски, выцветающая фотография.
География: лишний день моряка --- настоящий поворот кольца Земли.
Биология: подпитка, обмен, сон. Информатика: код с
самопроверкой, канал связи, два дома памяти. Обществознание:
прибавка связи, потолок пирамид. Физика: часы, сцена, частицы-узоры.
Книга не добавляет предмет --- она сшивает существующие.

\begin{askbox}[Подумай сам --- для взрослого]
Чему из этой книги ты сам сопротивляешься сильнее всего? Найди
пометку у этого утверждения. Если там \stT{} --- проверь
доказательство в корпусе; если \stH{} или \stI{} --- твоё
сопротивление законно и даже полезно. Так работает честное знание.
\end{askbox}

\clearpage
% ============================= МОСТЫ В КОРПУС =============================
\section*{Мосты в корпус}
\addcontentsline{toc}{section}{Мосты в корпус}

Для читателя, который хочет пройти от образа до доказательства.
Каждое несущее утверждение этой книги живёт в корпусе УГМ
(\holonsh) как теорема с доказательством, машинная проверка с
напечатанными числами или поимённо помеченный постулат; реестр
статусов там нумерует их все. Вот план дверей.

\begin{center}\small
\begin{tabular}{@{}p{0.36\textwidth}p{0.28\textwidth}p{0.28\textwidth}@{}}
\toprule
\textbf{Глава --- утверждение} & \textbf{Крыло корпуса} & \textbf{Ключевые записи} \\
\midrule
\ref{sec:fabric}--\ref{sec:seven}: семёрка вынуждена; ткань сама
себя проверяет & основания; доказательства минимальности & теорема
7-минимальности; $G_2$-жёсткость (T-42) \\[2pt]
\ref{sec:numbers}: лестница $1{\to}2{\to}4{\to}8$ обрывается;
сплошная прямая из уробороса & основания (миры чисел) & теорема о
алгебрах с делением; полнота через неподвижную точку \\[2pt]
\ref{sec:time}: время как взаимная согласованность & физическое
крыло & часы Пейджа--Вуттерса \\[2pt]
\ref{sec:source}: Источник дважды единствен и неустойчив &
основания; космогенезис & теоремы единственности и неустойчивости;
машинное интегрирование ската \\[2pt]
\ref{sec:space}: три оси; стрела $(1{,}3)$ & физическое крыло &
повороты, щадящие часы; теорема о подписи \\[2pt]
\ref{sec:life}: порог «вдвое собраннее хаоса»; свеча; прожиточный
ваттаж & динамика ядра; термодинамика & $P_{\text{crit}}=2/7$;
теоремы диссипации \\[2pt]
\ref{sec:window}: четыре стрелки окна & сознание: состояния &
$R\geq 1/3$ (T-45); $\Phi\geq 1$ (T-129); два тона \\[2pt]
\ref{sec:tower}: три этажа по $\times\tfrac13$; перечитывание
достраивает мысль & сознание: иерархия; самонаблюдение & башня
глубины; ретро-пополнение (C36) \\[2pt]
\ref{sec:qualia}: краски $=$ повороты колец; декодер; разрыв
нулевой ковариации; выцветание; паспорт & сознание: феноменология
--- механизм квалиа & T-300; T-301; T-302; выцветание $5\gamma/21$
(T-59); ёмкость (T-311) \\[2pt]
\ref{sec:freedom}--\ref{sec:golden}: двери свободы; водосборы,
смелость, часы опасности & операторы ядра; прикладные атласы &
ядро свободы; машинный рельеф и жизни \\[2pt]
\ref{sec:language}: $2{,}8$ бита за шаг; словам нужен правильный
глаз & феноменология: язык качества & теорема о канале; остенсивные
испытания \\[2pt]
\ref{sec:memory}: стол на семь мест, тающий втрое & сознание:
память; иерархия & закон трети внимания \\[2pt]
\ref{sec:others}--\ref{sec:worlds}: прибавка связи $\geq 0$;
потолок трёх этажей; эмбриология краски & структура ядра; иерархия;
феноменология & супераддитивность; потолок композиции; краска
имени \\[2pt]
\ref{sec:health}: панель стресса по голосам & прикладное крыло:
здоровье & эквивалентность стресса и жизнеспособности \\[2pt]
\ref{sec:death}: без Большого разрыва; бесцветное равновесие &
физика; феноменология & теорема исключения; следствия
выцветания \\[2pt]
\ref{sec:observer}: у фактов есть адреса & категорные основания &
сайт-релятивизация \\[2pt]
Приложение: колесо --- часы, комбинаторика,
рассказчик & исследовательское крыло & отчёты машинной
реконструкции \\[2pt]
\ref{sec:check}: столы ремёсел; четыре проверки; реестр & справка:
реестр фальсифицируемости & $\delta_{CP}$; почерк Фано в шуме;
излом $1/4$ пробуждения; две частицы \\
\bottomrule
\end{tabular}
\end{center}

\clearpage
% ============================= РЕЕСТР ПРОВЕРОК =============================
\section*{Реестр проверок}
\addcontentsline{toc}{section}{Реестр проверок}

Четырнадцать сводных предсказаний; у каждого --- полоса, нарушение
которой убивает соответствующее утверждение на его пометке.
Вердикты на 2026 год: \textbf{пройдено} --- измеренное значение
лежит внутри полосы; \emph{согласуется} --- не исключено, но за
пределами нынешней чувствительности; \emph{частично} --- косвенная
или калибровочная поддержка; \emph{не испытано} --- ни один прибор
ещё не дотянулся до полосы. Ни одно не провалено; ни одно не
списано втихую.

\begin{center}\small
\begin{tabular}{@{}p{0.44\textwidth}p{0.20\textwidth}p{0.26\textwidth}@{}}
\toprule
\textbf{Предсказание} & \textbf{Где проверять} & \textbf{Вердикт (2026)} \\
\midrule
Фаза CP $\delta_{CP}\approx 64{,}5^\circ\pm5^\circ$ \stH &
коллайдеры & \textbf{пройдено}: $64{,}6^\circ\pm2{,}8^\circ$ \\[2pt]
Угол Кабиббо $\approx 13^\circ$ \stH & каонные эксперименты &
\textbf{пройдено}: $12{,}96^\circ$ \\[2pt]
Электрический диполь нейтрона --- в точности ноль \stT & дипольные
поиски & \textbf{пройдено}: совместимо с $0$ \\[2pt]
Ровно одна тяжёлая связь порядка единицы \stT & коллайдеры &
\textbf{пройдено}: единственна, $\approx 0{,}94$ \\[2pt]
Две суммы всякого живого состояния в своих полосах \stT & любое
реконструированное состояние & \textbf{пройдено}: ни одного
контрпримера на $20\,000$ \\[2pt]
Без Большого разрыва; дрейф тёмной энергии предсказанной формы
\stT\,\stC & обзоры неба & согласуется \\[2pt]
Время жизни протона $\sim 6{,}7\times10^{37}$ лет \stH & подземные
детекторы & согласуется (выше нынешней досягаемости) \\[2pt]
Порог пробуждения на $2/7$; клиническая мера трётся о $0{,}31$
\stC & анестезиология & частично \\[2pt]
Шесть--двенадцать внутренних потоков покоящегося мозга \stH &
нейровизуализация & частично (сообщают $7$--$17$) \\[2pt]
Блочная структура Фано в когерентностях мозга \stT &
нейровизуализация & не испытано \\[2pt]
Внутритроечная скрытость ниже межтроечной \stH &
нейровизуализация & не испытано \\[2pt]
Почерк семёрки в скоростях распада (четырнадцать точных правил
сумм) \stT\,\stC & квантовая томография & не испытано \\[2pt]
Выделенный космический масштаб $\sim 160$ парсеков \stT &
крупномасштабные обзоры & не испытано \\[2pt]
Отклонение самосвязи Хиггса $10^{-2}$--$10^{-3}$ \stH & будущие
коллайдеры & не испытано \\
\bottomrule
\end{tabular}
\end{center}

\vfill
\begin{factbox}[Куда идти дальше]
Всё, что здесь рассказано образами, в корпусе УГМ (\holonsh) записано
теоремами с доказательствами, машинными проверками и реестром
статусов: космогенезис --- в разделе о происхождении вселенной;
окно, башня и краски --- в разделах о сознании; частицы --- в
физическом крыле; проверки --- в реестре предсказаний; а прикладное
крыло --- когерентная кибернетика --- несёт ту же семёрку в
протоколы, приборы и инженерию. Эта книга --- дверь; дом --- там.
\end{factbox}

% ============================= КОЛОФОН =============================
\clearpage
\thispagestyle{empty}
\vspace*{\fill}
\begin{center}\small\color{plosgray}
\begin{minipage}{0.78\textwidth}\centering
{\itshape Ткань бытия}\\[3pt]
Издание первое --- август 2026\\[12pt]
Набрано Xe\LaTeX{}ом гарнитурой CMU;\\
фигуры нарисованы в Ti\emph{k}Z по единой фирменной палитре:\\
семь цветов голосов, почти белые фоны, волосяные линейки.\\[8pt]
Каждое утверждение несёт одну из пяти пометок --- \stT\ \stC\ \stH\ \stP\ \stI\ ---\\
сжатых отпечатков одного тройного обязательства: построено, проверено, работает, ---\\
сверенных с реестром статусов корпуса УГМ (\holonsh).\\
Гейт сборки не допускает ни одной ошибки; каждая фигура вычитана глазами.\\[14pt]
\begin{tikzpicture}[scale=0.45]
  \voicerose{0}{0}{1.15}{0.1}
\end{tikzpicture}
\end{minipage}
\end{center}
\vspace*{\fill}


\clearpage
% ============================= ПРИЛОЖЕНИЕ: СТАРОЕ КОЛЕСО =============================
\section*{Приложение. Старое колесо под оптикой}
\addcontentsline{toc}{section}{Приложение. Старое колесо под оптикой}

\noindent\emph{Этот разбор стоит вне главной дороги книги:
рабочий пример оптики, приложенной к до-научной системе
самоописания. Он включён по одной причине --- на трудном случае
он показывает, как пометки честности отделяют прибор от
рассказа.}\medskip

\section{Кто смотрел на мир, пока смотреть было некому}
\label{sec:observer}

Мы дошли до края личной истории --- и упёрлись в вопрос обо всей
истории сразу.
Вопрос, который честно мучает всех, кто слышал слово «наблюдатель»
в физике: если сознание появляется поздно (глава~\ref{sec:source}),
то \emph{кто наблюдал вселенную первые миллиарды лет}? Кто
«схлопывал» её возможности в факты?

Ответ УГМ снимает вопрос, а не отвечает на него силой:
\textbf{фактам не нужен свидетель --- фактам нужен адрес} \stT.
В устройстве теории истина живёт не в чьей-то голове, а в самой
ткани: каждое «так» и «не так» приписано месту в структуре мира,
как высота приписана точке ландшафта. Гора была высокой до
альпинистов; «высота» не ждала, пока её измерят. Наблюдение ---
не акт творения факта, а \emph{поздний способ узнавать} факты,
доступный узорам из окна сознания.

Поэтому доответ на детский вопрос звучит так. До порога
жизнеспособности мир не «болтался неопределённым» и не «ждал
взгляда»: он определённо был --- безадресатно, как музыка в пустом
зале определённо звучит. А первые наблюдатели не застали мир
врасплох --- они \emph{в нём проснулись}, как просыпается слушатель,
который сам есть часть музыки. Мир не нуждался в свидетелях, чтобы
быть; но, по-видимому, он таков, что свидетели в нём неизбежно
заводятся (глава~\ref{sec:source}) --- и вот это, пожалуй, самое
удивительное из его свойств.

\begin{figure}[H]
\centering
\begin{tikzpicture}[font=\small]
  \fill[plosgray!18] (-4.6,0) -- (-2.2,2.5) -- (-0.6,0.9) -- (0.9,2.0) -- (3.6,0) -- cycle;
  \draw[plosgray!60, thick] (-4.6,0) -- (-2.2,2.5) -- (-0.6,0.9) -- (0.9,2.0) -- (3.6,0);
  \fill[softsky] (-2.55,2.13) -- (-2.2,2.5) -- (-1.9,2.18) -- cycle;
  \draw[fact, thick] (-2.2,2.5) -- (-2.2,3.05);
  \node[draw=fact, fill=softsky, rounded corners, font=\scriptsize, inner sep=2.5pt]
    at (-2.2,3.35) {высота: есть};
  \node[font=\scriptsize\itshape, text=plosgray, align=center] at (2.15,3.2)
    {альпинистов ещё нет ---\\ у факта есть адрес, а не зритель};
  \draw[plosgray!45] (-5.2,0) -- (4.2,0);
\end{tikzpicture}
\caption{Гора была высокой до всякого взгляда: истина в УГМ приписана
месту в ткани, как высота --- точке ландшафта. Наблюдение --- поздний
способ \emph{узнавать} факты, а не способ их создавать.}
\label{fig:mountain}
\end{figure}

Есть и последняя глубина, ради которой стоило писать
главу~\ref{sec:numbers}. Мир, моделирующий сам себя, держит себя за
хвост --- и требование, чтобы этот захват был возможен всегда,
оказалось не поэзией, а несущей деталью: оно вынуждает сплошную
числовую прямую, а с ней и непрерывное время \stT. Самонаблюдение ---
не поздняя прихоть мира. Сама \emph{возможность} смотреть на себя
вшита в счёт, которым мир записан, --- задолго до того, как в нём
проснулся первый смотрящий.

\begin{factbox}[Если точно]
В основании УГМ факты --- «секции» структуры, индексированные её же
внутренними адресами (сайт-релятивизация \stT); никакого внешнего
регистратора конструкция не требует и не допускает. Знаменитые
парадоксы наблюдателя разбираются именно этим ходом --- теория
предъявляет его как свой «четвёртый путь» в современном споре об
объективности \stT. Сознание при этом остаётся особенным --- но как
\emph{поздний житель} мира, а не как его фундамент.
\end{factbox}


\end{document}
